\documentclass[11pt,a4paper]{article}
% ============================================================================
%  TFPT consolidated document --- unified preamble (superset of the merged notes)
% ============================================================================
\usepackage[T1]{fontenc}
\usepackage[utf8]{inputenc}
\usepackage[english]{babel}
\usepackage{lmodern}
\usepackage[a4paper,margin=2.3cm]{geometry}
\usepackage{amsmath,amssymb,amsthm,mathtools}
\usepackage{booktabs,array,longtable,tabularx,multirow}
\usepackage{enumitem}
\usepackage{xcolor}
\usepackage{tikz}
\usetikzlibrary{arrows.meta,positioning,calc,fit,backgrounds,shapes.geometric,shapes.misc,decorations.pathmorphing,decorations.pathreplacing,patterns,angles,quotes,matrix}
\usepackage{microtype}
\usepackage[most]{tcolorbox}
\usepackage[hidelinks,colorlinks=true,linkcolor=blue!55!black,citecolor=blue!55!black,urlcolor=blue!55!black]{hyperref}
\emergencystretch=3.5em

% ---- colours ----
% ---------------------------------------------------------------------
%  tfpt_style.tex -- THE shared visual style for every TFPT document.
%  Single source for colours, status markers, marker key, the \veri
%  script reference, the tcolorbox families, and the running
%  header/footer.  \input AFTER xcolor + tcolorbox[most] + hyperref are
%  loaded (every paper loads those).  Each document sets its short
%  running title before \begin{document}, e.g.
%        \def\TFPTdoctitle{Architecture and the E8 Compiler}
%  ---------------------------------------------------------------------
\usepackage{fancyhdr}
\usepackage{etoolbox}
\usepackage{tikz}
\usetikzlibrary{positioning,arrows.meta,fit,backgrounds,calc,shapes.geometric,shapes.misc}

% ---- single-source author / version / automatic date stamp ----
% ---------------------------------------------------------------------
%  tex-artefacts/version.tex -- single source for author list, version
%  and the automatic title-page date stamp of every TFPT document.
%  \input by tex-artefacts/tfpt_style.tex, so every paper that loads the
%  shared style automatically gets:
%     \TFPTauthors   -- the author list (use as \author{\TFPTauthors})
%     \TFPTversion   -- curated semantic version (bump per release; mirror
%                       the bump in changelog.tex)
%     \TFPTrev       -- auto build revision = git commit count, injected by
%                       build.sh into tex-artefacts/version-auto.tex
%                       (gitignored); falls back to "dev" for standalone runs
%     \TFPTdatestamp -- the title-page date line: \today + version + revision
%  build.sh also writes tex-artefacts/release-version.txt (gitignored):
%     "{TFPTversion} (rev {git commit count})" for Zenodo + release tooling.
%  Bump \TFPTversion here ONCE; it then updates in all papers at the next build.
% ---------------------------------------------------------------------
\def\TFPTauthors{Stefan Hamann \and Alessandro Rizzo}
\def\TFPTversion{5.3}
\def\TFPTrev{dev}
% build.sh regenerates the line below with the live git commit count:
\InputIfFileExists{tex-artefacts/version-auto.tex}{}{}
\def\TFPTdatestamp{\today\,\textperiodcentered\,v\TFPTversion\,(rev~\TFPTrev)}


% ---- colours (canonical TFPT palette) ----
\definecolor{tfblue}{HTML}{1F4E79}
\definecolor{tfgreen}{HTML}{1E7B45}
\definecolor{tfred}{HTML}{9B2226}
\definecolor{tfgold}{HTML}{B8860B}
\definecolor{tfgray}{HTML}{555555}

% ---- status markers (simplified 4-class display system) ----
% Reader-facing markers collapse to FOUR classes.  The fine-grained per-claim
% type (Axiom / Formal / Lattice / Numerical / Identity / Physical) stays in
% verification/status_ledger.csv -- the single source of truth -- so no fidelity
% is lost; the papers just show the coarse class so a reader is not overwhelmed.
%   [E] exact / machine-proven  (identity, Lie/lattice, formalised, numerical FP)
%   [C] conditional             (physical / bridge / readout, under named hypotheses)
%   [O] open / axiom            (declared input or genuine gap)
%   [X] kill test               (falsifiable threshold)
\newcommand{\mE}{{\small\textcolor{tfgreen!70!black}{\textbf{[E]}}}}
\newcommand{\mC}{{\small\textcolor{tfgold!78!black}{\textbf{[C]}}}}
\newcommand{\mO}{{\small\textcolor{tfred!80!black}{\textbf{[O]}}}}
\newcommand{\mX}{{\small\textcolor{tfred!58!black}{\textbf{[X]}}}}
% backward-compatible aliases: the old fine-grained names now render their class,
% so any not-yet-rewritten call site (or the standalone changelog) still compiles.
\newcommand{\tI}{\mE}\newcommand{\tL}{\mE}\newcommand{\tF}{\mE}\newcommand{\tN}{\mE}
\newcommand{\tIalg}{\mE}\newcommand{\tInum}{\mE}
\newcommand{\tP}{\mC}\newcommand{\tB}{\mC}\newcommand{\tR}{\mC}
\newcommand{\tA}{\mO}
\newcommand{\tX}{\mX}

% compact marker legend, placed under a marker-bearing table
\newcommand{\markerkey}{\par\vspace{1.5pt}\noindent{\scriptsize\itshape
  \textcolor{black!58}{Marker key:\ \mE~exact / machine-proven\,$\cdot$\,%
  \mC~conditional (holds under named hypotheses)\,$\cdot$\,%
  \mO~open / axiom\,$\cdot$\,\mX~falsifiable kill test.\ \
  Per-claim fine type in \texttt{verification/status\_ledger.csv}.}}\par\vspace{2pt}}

% horizontal badge bar (place once near the front of a document)
\newcommand{\TFPTbadges}{\par\noindent
  \fcolorbox{tfgreen!55!black}{tfgreen!8}{\scriptsize\mE\,exact / machine-proven}\hfill
  \fcolorbox{tfgold!70!black}{tfgold!10}{\scriptsize\mC\,conditional}\hfill
  \fcolorbox{tfred!75!black}{tfred!8}{\scriptsize\mO\,open / axiom}\hfill
  \fcolorbox{tfred!60!black}{tfred!6}{\scriptsize\mX\,kill test}\par\vspace{3pt}}

% inline reference to the script that machine-checks a claim (see verification/)
\newcommand{\veri}[1]{\,{\footnotesize\textcolor{tfblue!70!black}{[\,\texttt{verification/#1}\,]}}}
% lightweight inline colour-mark for a bare verification-script token in running
% prose / tables, e.g. \vref{v108}, \vref{v49}/\vref{v71}.  Same blue as \veri so a
% reader instantly sees "this is a machine-checked module reference".
\newcommand{\vref}[1]{\textcolor{tfblue!72!black}{\textsf{#1}}}

% ---- tcolorbox families (one visual language across all documents) ----
% keybox  : the central claim / reading of a section (blue)
% winbox  : an established result / closure (green)
% gapbox  : an open gate / honest caveat / kill criterion (red)
% auditbox: a recorded audit fingerprint, not load-bearing (gold)
% navbox  : a light, neutral navigation / reference card (document map, etc.)
% Visual language (deliberately light): a faint tint, a thin coloured frame
% and a light title band with dark coloured title text -- NOT a heavy
% white-on-saturated-colour bar -- so a page never reads as a wall of boxes.
\newtcolorbox{keybox}[1][]{enhanced,breakable,boxrule=0.5pt,arc=1.5pt,
  colback=tfblue!4!white,colframe=tfblue!70!black,coltitle=tfblue!88!black,colbacktitle=tfblue!12!white,
  fonttitle=\bfseries\small,fontupper=\small,title={#1},left=2mm,right=2mm,top=1mm,bottom=1mm}
% non-breaking variant (front-matter cards that should never split across a page)
\newtcolorbox{keyboxnb}[1][]{enhanced,breakable=false,boxrule=0.5pt,arc=1.5pt,
  colback=tfblue!4!white,colframe=tfblue!70!black,coltitle=tfblue!88!black,colbacktitle=tfblue!12!white,
  fonttitle=\bfseries\small,fontupper=\small,title={#1},left=2mm,right=2mm,top=1mm,bottom=1mm}
\newtcolorbox{winbox}[1][]{enhanced,breakable,boxrule=0.5pt,arc=1.5pt,
  colback=tfgreen!5!white,colframe=tfgreen!65!black,coltitle=tfgreen!50!black,colbacktitle=tfgreen!14!white,
  fonttitle=\bfseries\small,fontupper=\small,title={#1},left=2mm,right=2mm,top=1mm,bottom=1mm}
\newtcolorbox{gapbox}[1][]{enhanced,breakable,boxrule=0.5pt,arc=1.5pt,
  colback=tfred!4!white,colframe=tfred!70!black,coltitle=tfred!75!black,colbacktitle=tfred!12!white,
  fonttitle=\bfseries\small,fontupper=\small,title={#1},left=2mm,right=2mm,top=1mm,bottom=1mm}
% non-breaking variants (callouts that must stay whole on one page, embedded in text)
\newtcolorbox{winboxnb}[1][]{enhanced,breakable=false,boxrule=0.5pt,arc=1.5pt,
  colback=tfgreen!5!white,colframe=tfgreen!65!black,coltitle=tfgreen!50!black,colbacktitle=tfgreen!14!white,
  fonttitle=\bfseries\small,fontupper=\small,title={#1},left=2mm,right=2mm,top=1mm,bottom=1mm}
\newtcolorbox{gapboxnb}[1][]{enhanced,breakable=false,boxrule=0.5pt,arc=1.5pt,
  colback=tfred!4!white,colframe=tfred!70!black,coltitle=tfred!75!black,colbacktitle=tfred!12!white,
  fonttitle=\bfseries\small,fontupper=\small,title={#1},left=2mm,right=2mm,top=1mm,bottom=1mm}
\newtcolorbox{auditbox}[1][]{enhanced,breakable,boxrule=0.5pt,arc=1.5pt,
  colback=tfgold!7!white,colframe=tfgold!70!black,coltitle=tfgold!45!black,colbacktitle=tfgold!16!white,
  fonttitle=\bfseries\small,fontupper=\small,title={#1},left=2mm,right=2mm,top=1mm,bottom=1mm}
\newtcolorbox{navbox}[1][]{enhanced,breakable,boxrule=0.5pt,arc=1.5pt,
  colback=tfgray!3!white,colframe=tfgray!55,coltitle=tfgray!30!black,colbacktitle=tfgray!12!white,
  fonttitle=\bfseries\small,fontupper=\small,title={#1},left=2mm,right=2mm,top=1mm,bottom=1mm}
% back-compatibility aliases (older box names map onto the canonical types)
\newtcolorbox{audbox}[1][]{enhanced,breakable,boxrule=0.5pt,arc=1.5pt,
  colback=tfgold!7!white,colframe=tfgold!70!black,coltitle=tfgold!45!black,colbacktitle=tfgold!16!white,
  fonttitle=\bfseries\small,fontupper=\small,title={#1},left=2mm,right=2mm,top=1mm,bottom=1mm}
\newtcolorbox{resbox}[1][]{enhanced,breakable,boxrule=0.5pt,arc=1.5pt,
  colback=tfgreen!5!white,colframe=tfgreen!55!black,coltitle=tfgreen!50!black,colbacktitle=tfgreen!12!white,
  fonttitle=\bfseries\small,fontupper=\small,title={#1},left=2mm,right=2mm,top=1mm,bottom=1mm}
\newtcolorbox{roadbox}[1][]{enhanced,breakable,boxrule=0.5pt,arc=1.5pt,
  colback=tfgreen!5!white,colframe=tfgreen!55!black,coltitle=tfgreen!50!black,colbacktitle=tfgreen!12!white,
  fonttitle=\bfseries\small,fontupper=\small,title={#1},left=2mm,right=2mm,top=1mm,bottom=1mm}
\newtcolorbox{intbox}[1][]{enhanced,breakable,boxrule=0.5pt,arc=1.5pt,
  colback=tfgold!7!white,colframe=tfgold!70!black,coltitle=tfgold!45!black,colbacktitle=tfgold!16!white,
  fonttitle=\bfseries\small,fontupper=\small,title={#1},left=2mm,right=2mm,top=1mm,bottom=1mm}

% ---- colour-marked display formula (load-bearing key equations) ----
% A highlighted formula line: a coloured left accent bar + faint tint, no full
% frame -- so a key equation is visually set apart from the running prose
% without becoming yet another box.  Use:  \begin{keyeq}$\displaystyle ...$\end{keyeq}
\newtcolorbox{keyeq}{blanker,breakable=false,halign=center,
  colback=tfgold!8!white,borderline west={2.4pt}{0pt}{tfgold!75!black},
  left=8pt,right=6pt,top=3pt,bottom=3pt,before skip=5pt,after skip=5pt}

% ---- running header / footer (consistent across all documents) ----
\providecommand{\TFPTdoctitle}{}
\setlength{\headheight}{14pt}
\pagestyle{fancy}
\fancyhf{}
\fancyhead[L]{\footnotesize\textcolor{tfgray}{\textit{TFPT}\,\textbullet\,\TFPTdoctitle}}
\fancyhead[R]{\footnotesize\textcolor{tfgray}{\thepage}}
\fancyfoot[L]{\scriptsize\textcolor{tfgray}{v\TFPTversion\,\textperiodcentered\,rev~\TFPTrev}}
\renewcommand{\headrulewidth}{0.3pt}
\renewcommand{\footrulewidth}{0pt}
\fancypagestyle{plain}{\fancyhf{}\renewcommand{\headrulewidth}{0pt}%
  \fancyfoot[C]{\footnotesize\textcolor{tfgray}{\thepage}}%
  \fancyfoot[L]{\scriptsize\textcolor{tfgray}{v\TFPTversion\,\textperiodcentered\,rev~\TFPTrev}}}

% ---- shared figure macros (claim stack, anchor cockpit, E8 glue, 2/3 motif, K5) ----
%  (this fragment lives in tex-artefacts/ and is \input by documents compiled from
%   the repository root, so the path is given relative to that root.)
% ---------------------------------------------------------------------
%  tfpt_figures.tex -- shared, reusable TikZ figure macros for every
%  TFPT document.  \input by tfpt_style.tex (which loads tikz + the
%  needed libraries + etoolbox).  Each macro draws a self-contained
%  centred figure; nothing is drawn until the macro is called.
%
%    \TFPTclaimstack[zone]  -- the architecture claim stack (B3); the
%                              optional zone in {axioms,anchor,compiler,
%                              sm,bridges,frontier} is highlighted.
%    \TFPTanchorcockpit      -- the anchor a=(1,1,2) cockpit (B4)
%    \TFPTeightglue          -- the D5 (+) A3 +mu4 => E8 glue (B5)
%    \TFPTmotif              -- the 2/3 motif map (B6)
%    \TFPTactionkfive        -- the K5 action-lattice mnemonic (B7)
% ---------------------------------------------------------------------

% ===== B3: the architecture claim stack =========================
\newcommand{\tfhl}[1]{\ifdefstring{\tfcur}{#1}{cbhi}{cbnl}}
\newcommand{\TFPTclaimstack}[1][none]{%
\def\tfcur{#1}%
\begin{center}
\begin{tikzpicture}[font=\footnotesize,
  cb/.style={rounded corners=2pt,inner sep=4pt,text width=118mm,align=center,minimum height=8.5mm},
  cbnl/.style={line width=0.5pt}, cbhi/.style={line width=1.5pt},
  ar/.style={-{Stealth[length=2.2mm]},tfgray,semithick}]
  \node[cb,fill=tfred!8,draw=tfred!75!black,\tfhl{axioms}] (axioms)
    {\textbf{P1}\,$c_3=\tfrac1{8\pi}$ \quad \textbf{P2}\,$g_{\mathrm{car}}=5$ \ \textcolor{tfgray}{\itshape(two inputs, [O] axioms)}};
  \node[cb,fill=tfblue!8,draw=tfblue,\tfhl{anchor},below=3.5mm of axioms] (anchor)
    {Anchor $a=(1,1,2)$:\ $e_1,e_2,e_3=4,5,2$ \ \textcolor{tfgray}{\itshape([E] exact)}};
  \node[cb,fill=tfgreen!10,draw=tfgreen!70!black,\tfhl{compiler},below=3.5mm of anchor] (compiler)
    {Discrete compiler:\ $D_5\oplus A_3+\mu_4\Rightarrow E_8$,\ $h=30$,\ $\varphi_0$ \ \textcolor{tfgray}{\itshape([E] exact)}};
  \node[cb,fill=tfgreen!7,draw=tfgreen!60!black,\tfhl{sm},below=3.5mm of compiler] (sm)
    {SM skeleton:\ three families, hypercharges, flavor matrix $R$ \ \textcolor{tfgray}{\itshape([E])}};
  \node[cb,fill=tfgold!12,draw=tfgold!80!black,\tfhl{bridges},below=3.5mm of sm] (bridges)
    {Physical bridges:\ $v_{\mathrm{geo}}$,\ $G_{\mathrm{net}}$,\ $F_{\mathrm{transfer}}$ \ \textcolor{tfgray}{\itshape([C] bridges)}};
  \node[cb,fill=tfgray!12,draw=tfgray!70!black,\tfhl{frontier},below=3.5mm of bridges] (frontier)
    {Frontier readouts:\ $\Lambda$, inflation, $\eta_B$, dark matter, QG \ \textcolor{tfgray}{\itshape([C]/[O])}};
  \draw[ar] (axioms)--(anchor); \draw[ar] (anchor)--(compiler);
  \draw[ar] (compiler)--(sm);   \draw[ar] (sm)--(bridges); \draw[ar] (bridges)--(frontier);
\end{tikzpicture}
\end{center}}

% ===== B4: the anchor cockpit ====================================
\newcommand{\TFPTanchorcockpit}{%
\begin{center}
\begin{tikzpicture}[font=\footnotesize,
  pan/.style={rounded corners=2pt,draw=tfblue!60,fill=tfblue!4,inner sep=5pt,
    text width=46mm,align=left},
  hd/.style={font=\bfseries\small,tfblue}]
  \node[pan] (e) {\textbf{elementary} $e_k(a)$\\[2pt]
     $e_1=4\to|\mu_4|$\\ $e_2=5\to g_{\mathrm{car}}$\\ $e_3=2\to|\mathbb Z_2|$\\[2pt]
     $c_3=\tfrac1{2e_1\pi}=\tfrac1{8\pi}$};
  \node[pan,right=4mm of e] (p) {\textbf{power sums} $p_n=2+2^n$\\[2pt]
     $p_0,p_1,p_2,p_3=3,4,6,10$\\ $p_1p_2p_3=240=|R(E_8)|$\\ $p_4-p_3=8=\operatorname{rank}E_8$\\
     $\Rightarrow\dim E_8=248$};
  \node[pan,right=4mm of p] (m) {\textbf{derived motifs}\\[2pt]
     $4,5,8,16,25,$\\ $30,41,120,240$\\[2pt]
     \textcolor{tfgray}{\itshape one anchor,}\\ \textcolor{tfgray}{\itshape many projections}};
  \node[hd,above=1mm of p.north] {$a=(1,1,2)$ \ \textcolor{tfgray}{\footnotesize\itshape[E]}};
\end{tikzpicture}
\end{center}}

% ===== B5: the E8 glue ===========================================
\newcommand{\TFPTeightglue}{%
\begin{center}
\begin{tikzpicture}[font=\footnotesize,
  d/.style={rounded corners=2pt,draw=tfblue,fill=tfblue!6,inner sep=4pt,align=center,text width=36mm},
  mid/.style={rounded corners=2pt,draw=tfgold!80!black,fill=tfgold!12,inner sep=4pt,align=center,text width=52mm},
  e8/.style={rounded corners=2pt,draw=tfgreen!70!black,fill=tfgreen!10,inner sep=4pt,align=center,text width=52mm,font=\bfseries},
  rt/.style={rounded corners=2pt,draw=tfgreen!55!black,fill=tfgreen!6,inner sep=3pt,align=center,text width=52mm},
  side/.style={rounded corners=2pt,draw=tfgray!55,fill=tfgray!6,inner sep=4pt,align=left,text width=34mm,font=\scriptsize},
  ar/.style={-{Stealth[length=2mm]},tfgray,semithick}]
  \node[d] (d5) {$D_5$ discriminant\\ $\mathbb Z_4$, $q=5/4$};
  \node[d,right=10mm of d5] (a3) {$A_3$ discriminant\\ $\mathbb Z_4$, $q=3/4$};
  \node[mid,below=5mm of $(d5)!0.5!(a3)$] (glue) {$\mu_4$ anti-isometry\ ($q$-sum $=2$)};
  \node[e8,below=4mm of glue] (e8) {even unimodular $E_8$ lattice};
  \node[rt,below=4mm of e8] (rt) {$240$ roots,\ rank $8$,\ $\dim=248$};
  \node[side,right=8mm of glue] (why) {\textbf{why unique?}\\ rank $8$\\ even\\ unimodular\\ positive definite};
  \draw[ar] (d5)--(glue); \draw[ar] (a3)--(glue);
  \draw[ar] (glue)--(e8); \draw[ar] (e8)--(rt);
\end{tikzpicture}\\[2pt]
{\scriptsize\itshape $D_5\oplus A_3+\mu_4\Rightarrow E_8$ --- the unique even unimodular rank-$8$ glue \ [E] exact (lattice).}
\end{center}}

% ===== B5b: the double pushout (one glue, two categories) ========
\newcommand{\TFPTdoublepushout}{%
\begin{center}
\begin{tikzpicture}[font=\footnotesize,>={Stealth[length=2mm]},
  lat/.style={rounded corners=2pt,draw=tfblue,fill=tfblue!6,inner sep=5pt,align=center,text width=34mm},
  net/.style={rounded corners=2pt,draw=tfblue!70!black,fill=tfblue!4,inner sep=5pt,align=center,text width=34mm},
  e8l/.style={rounded corners=2pt,draw=tfgold!80!black,fill=tfgold!12,inner sep=5pt,align=center,text width=34mm,font=\bfseries},
  e8n/.style={rounded corners=2pt,draw=tfgreen!70!black,fill=tfgreen!10,inner sep=5pt,align=center,text width=34mm,font=\bfseries},
  ar/.style={->,tfgray,semithick},
  lb/.style={font=\scriptsize,text=black,align=center,fill=white,inner sep=1.5pt}]
  \node[lat] (LL) at (0,2.5)   {$D_5\oplus A_3$\\[1pt]\scriptsize root lattices,\ $\mathrm{disc}=\Z_4$};
  \node[e8l] (LR) at (6.6,2.5) {$E_8$\\[1pt]\scriptsize even unimodular,\ $\det 1$};
  \node[net] (BL) at (0,0)     {$(D_5)_1\otimes(A_3)_1$\\[1pt]\scriptsize chiral nets,\ $c=8$};
  \node[e8n] (BR) at (6.6,0)   {$(E_8)_1$\\[1pt]\scriptsize holomorphic,\ $\mu$-index $1$};
  \draw[ar] (LL)--node[lb]{$+\,\mu_4$ glue}(LR);
  \draw[ar] (BL)--node[lb]{$\rtimes\,\mu_4$ extension}(BR);
  \draw[ar] (LL)--node[lb,left=-1pt]{level-$1$}(BL);
  \draw[ar] (LR)--node[lb,right=-1pt]{level-$1$}(BR);
\end{tikzpicture}\\[2pt]
{\scriptsize\itshape One operation, two categories: the order-$|\mu_4|{=}4$ glue is the lattice pushout
$D_5\oplus A_3+\mu_4\Rightarrow E_8$ \emph{and} the simple-current / anyon-condensation extension
$(D_5)_1\otimes(A_3)_1\rtimes\mu_4\cong(E_8)_1$; the square commutes. \ [E] exact.}
\end{center}}

% ===== B6: the 2/3 motif map =====================================
\newcommand{\TFPTmotif}{%
\begin{center}
\begin{tikzpicture}[font=\footnotesize,
  ex/.style={rounded corners=2pt,draw=tfgreen!70!black,fill=tfgreen!8,inner sep=4pt,align=center,text width=92mm},
  br/.style={rounded corners=2pt,draw=tfgold!80!black,fill=tfgold!12,inner sep=4pt,align=center,text width=92mm},
  ar/.style={-{Stealth[length=2mm]},tfgray,semithick}]
  \node[ex] (a) {$|\mathbb Z_2|/N_{\mathrm{fam}}=2/3$ \ \textcolor{tfgray}{\itshape(anchor ratio)}};
  \node[ex,below=3mm of a] (b) {parabolic weights $\{0,\tfrac13,\tfrac23\}=\operatorname{spec}A_0^\star$};
  \node[ex,below=3mm of b] (c) {transfer spectrum $\{1,(2/3)^6,(1/3)^6\}$,\ gap $6\log\tfrac32$};
  \node[br,below=3mm of c] (d) {Koide target $2/3$ \ \textcolor{tfgray}{\itshape(source$\to$pole bridge [C])}};
  \node[br,below=3mm of d] (e) {Nariai / horizon entropy bound $2/3$ \ \textcolor{tfgray}{\itshape([C])}};
  \draw[ar] (a)--(b); \draw[ar] (b)--(c); \draw[ar] (c)--(d); \draw[ar] (d)--(e);
\end{tikzpicture}\\[2pt]
{\scriptsize\itshape Green: exact motif ([E]).\quad Gold: physical bridge ([C]) --- the same ratio, two epistemic types.}
\end{center}}

% ===== B7: the K5 action-lattice mnemonic ========================
%  Layout: complete graph K5 (pentagon, all 10 edges) on the LEFT, the
%  Pascal-row reading list on the RIGHT, vertically centred on the graph
%  so the two blocks never overlap.
\newcommand{\TFPTactionkfive}{%
\begin{center}
\begin{tikzpicture}[font=\scriptsize]
  % --- K5 vertices on a regular pentagon, centred at the origin ---
  \foreach \i in {1,...,5}{\coordinate (k\i) at ({90+72*(\i-1)}:12mm);}
  \begin{scope}[on background layer]
    \foreach \i in {1,...,5}{\foreach \j in {1,...,5}{\ifnum\i<\j
      \draw[tfgray!55,line width=0.5pt] (k\i)--(k\j);\fi}}
  \end{scope}
  \foreach \i in {1,...,5}{\fill[tfblue!70!black] (k\i) circle (1.8pt);}
  \node[tfgray,font=\scriptsize\itshape] at (0,-1.9) {$K_5$};
  % --- reading list, anchored to the right of the whole graph, centred ---
  \node[anchor=west,align=left,inner sep=0pt] at (1.9cm,0) {%
    $\tbinom51=5$\ \ vertices --- action $x^1$ (Hubble slot)\\[1pt]
    $\tbinom52=10$\ edges --- $\Lambda$ density pair $x^{10}$\\[1pt]
    $\tbinom53=10$\ triangles --- dual sector\\[1pt]
    $\tbinom54=5$\ \ complements --- inverse sector\\[1pt]
    $\tbinom55=1$\ \ determinant --- closure};
\end{tikzpicture}\\[2pt]
{\scriptsize\itshape $K_5$ on the five carrier slots; the action powers $x^{1},x^{5},x^{10}$ are its
$\tbinom5k$ faces \ \textcolor{tfgray}{(mnemonic; the physical scale map stays [C]).}}
\end{center}}



% ---- shared header/footer (version stamp) + 4-class markers come from tfpt_style ----
%  (no per-document footer override, so every TFPT document looks identical.)
\usepackage{needspace}
\usepackage{etoolbox}
% ---- guard every coloured box so it never breaks with an orphaned title line ----
\BeforeBeginEnvironment{keybox}{\needspace{8\baselineskip}}
\BeforeBeginEnvironment{winbox}{\needspace{8\baselineskip}}
\BeforeBeginEnvironment{gapbox}{\needspace{8\baselineskip}}
\BeforeBeginEnvironment{auditbox}{\needspace{8\baselineskip}}
\BeforeBeginEnvironment{audbox}{\needspace{8\baselineskip}}
\BeforeBeginEnvironment{resbox}{\needspace{8\baselineskip}}
\BeforeBeginEnvironment{roadbox}{\needspace{8\baselineskip}}

% ---- symbol macros (union; one definition each) ----
\newcommand{\TFPT}{\mathrm{TFPT}}
\newcommand{\MPl}{M_{\mathrm{Pl}}}
\newcommand{\Mpl}{M_{\mathrm{Pl}}}
\newcommand{\Mbar}{\bar M_{\mathrm{Pl}}}
\newcommand{\ellbar}{\bar\ell_{\mathrm{Pl}}}
\newcommand{\GN}{G_N}
\newcommand{\gcar}{g_{\mathrm{car}}}
\newcommand{\Nfam}{N_{\mathrm{fam}}}
\newcommand{\Oadm}{\Omega_{\mathrm{adm}}}
\newcommand{\cthree}{c_3}
\newcommand{\phiz}{\varphi_0^{\mathrm{ret}}}
\newcommand{\phibase}{\varphi_{\mathrm{base}}}
\newcommand{\phiseam}{\varphi_{\mathrm{seam}}}
\newcommand{\useed}{u}
\newcommand{\dtop}{\delta_{\mathrm{top}}}
\newcommand{\dtwo}{\delta_2}
\newcommand{\dph}{\delta_{\mathrm{ph}}}
\newcommand{\betarad}{\beta_{\mathrm{rad}}}
\newcommand{\lC}{\lambda_C}
\newcommand{\lamC}{\lambda_C}
\newcommand{\ainv}{\alpha^{-1}_\star}
\newcommand{\LamIR}{\Lambda_{\mathrm{IR}}}
\newcommand{\Lamgeom}{\Lambda_{\mathrm{geom}}}
\newcommand{\lamsig}{\lambda_\Sigma}
\newcommand{\AEW}{A_{\mathrm{EW}}}
\newcommand{\AL}{\mathcal A_\Lambda}
\newcommand{\AH}{\mathcal A_H}
\newcommand{\dimg}{\dim\mathfrak g_{\mathrm{SM}}}
\newcommand{\Dstart}{D_{\mathrm{start}}}
\newcommand{\kappaE}{\kappa_{E_8}}
\newcommand{\PP}{\mathbb{P}}
\newcommand{\Z}{\mathbb{Z}}
\newcommand{\onev}{\mathbf 1}
\newcommand{\lambdaY}{\lambda_Y}
\newcommand{\lamY}{\lambda_Y}
\newcommand{\vgeo}{v_{\mathrm{geo}}}
\newcommand{\pardeg}{\operatorname{pardeg}}

% ---- theorem environments (union) ----
\theoremstyle{definition}
\newtheorem{theorem}{Theorem}
\newtheorem{proposition}{Proposition}
\newtheorem{lemma}{Lemma}
\newtheorem{corollary}{Corollary}
\newtheorem{conjecture}{Conjecture}
\newtheorem{definition}{Definition}
\newtheorem*{remark}{Remark}
\newtheorem*{goal}{Project goal}

% ---- table column types ----
\newcolumntype{Y}{>{\raggedright\arraybackslash}X}
\newcolumntype{C}{>{\centering\arraybackslash}X}

% ---- list spacing ----
\setlist[itemize]{topsep=4pt,itemsep=2pt,parsep=0pt}
\setlist[enumerate]{topsep=4pt,itemsep=2pt,parsep=0pt}

% ---- TikZ styles used by the architecture diagrams ----
\tikzset{
  cbox/.style={rounded corners=2pt, draw=blue!55!black, fill=blue!6, line width=0.6pt, align=center, inner sep=4pt, font=\small},
  dbox/.style={rounded corners=2pt, draw=teal!60!black, fill=teal!8, line width=0.6pt, align=center, inner sep=4pt, font=\small},
  abox/.style={rounded corners=2pt, draw=orange!70!black, fill=orange!10, line width=0.6pt, align=center, inner sep=4pt, font=\small},
  ebox/.style={rounded corners=2pt, draw=purple!60!black, fill=purple!8, line width=0.8pt, align=center, inner sep=5pt, font=\small\bfseries},
  flow/.style={-{Stealth[length=2.4mm]}, line width=0.7pt, draw=black!65},
}
\renewcommand{\arraystretch}{1.15}

\title{\vspace{-1.2cm}\textbf{TFPT --- $E_8$ Audit, Cascade Bridge and Bootstrap}\\[2mm]\large The seven $E_8$ slices as an audit raster, the old cascade as the same even-integer spine, and the M\"obius self-consistency loop}
\author{\TFPTauthors}\date{\TFPTdatestamp}
\def\TFPTdoctitle{$E_8$ Audit \& Bootstrap}
\begin{document}
\maketitle
\vspace{-0.4cm}

\begin{keybox}[What this document is about]
$E_8$ as an \emph{audit container}, not a mystery: the seven maximal slices of $248$ as a falsification
raster (every load-bearing number must appear in at least one projection), the bridge showing the old
$E_8$ orbit cascade $D{=}60{-}2n$ is the same even-integer spine as the compiler, and the M\"obius bootstrap
in which $g_{\mathrm{car}}=5$ and the ``$8$'' in $c_3$ are overdetermined $E_8$-closure fixed points (only
$\pi$ irreducible).
\end{keybox}

\begin{gapbox}[What $E_8$ is \emph{not} (and why the known no-go results do not apply)]
$E_8$ here is the unimodular \emph{audit / compiler hull} that classifies the admissible discrete charge
and residue structures --- \textbf{not} an unbroken physical gauge group. The Standard Model is a
\emph{readout} after projection (not a direct $E_8$ gauge theory), the Lorentz signature appears only
after projection, and fermions are not ``everything in the adjoint''. Consequently the standard
objections to literal $E_8$ world-formulas do \emph{not} bite: Distler--Garibaldi (no chiral fermions /
representation obstructions for $E_8$ as a 4d gauge group) and Coleman--Mandula (no nontrivial mixing of
spacetime and internal symmetry) both constrain $E_8$ as a \emph{spacetime gauge symmetry}, which TFPT
does not claim. The defensible statement is: \emph{TFPT is the unique parabolic compilation of the anchor
$a=(1,1,2)$ into the $E_8$ audit hull} --- a discrete classifier, attacked on its own (combinatorial)
terms, not on gauge-theoretic ones.
\end{gapbox}

% ---------------------------------------------------------------------
%  tfpt_docset.tex -- THE canonical TFPT document-set box.
%  Single source of truth for the document list; \input by the
%  introduction and every paper. Uses only universally available
%  macros (keybox + plain LaTeX math) so it compiles in every doc.
%  Each document adds its own "You are reading ..." line directly
%  after the \input.  Numbering is aligned with the file names:
%  document N == tfpt_N_*, the introduction is the entry point, and
%  R/H/O are the appendix-style companions.
% ---------------------------------------------------------------------
\begin{navbox}[The TFPT document set --- what is treated where]
\textbf{Plain language:} TFPT (Topological Fixed-Point Theory) is a small discrete compiler. Two
inputs --- the seam constant $c_3=\tfrac1{8\pi}$ and the carrier rank $g_{\mathrm{car}}=5$ --- build
$D_5\oplus A_3+\mu_4\Rightarrow E_8$ and read off the Standard-Model skeleton, the constants and the
scale grammar. The series is \textbf{the introduction plus five numbered papers and three companions},
best read in order:
\begin{enumerate}[leftmargin=2.0em,itemsep=1pt,label=\textbf{\arabic*.}]
\setcounter{enumi}{-1}
\item \texttt{introduction} --- reading guide, compiler closure, predictions, the dependency DAG and the
      single proof ledger (\texttt{verification/status\_ledger.csv}). \emph{(entry point)}
\item \texttt{tfpt\_1\_architecture\_e8} --- the two axioms, the derivation map, the EM fixed point
      $\alpha^{-1}$, the $D_5\oplus A_3+\mu_4\Rightarrow E_8$ construction, the QBL theorem chain.
\item \texttt{tfpt\_2\_standard\_model} --- the SM in one $\varphi_0$-ladder formula, flavor from
      parabolic transport (sheet diamond, dual anchor, $Q_+$ cohomology), the worked closures.
\item \texttt{tfpt\_3\_e8\_audit\_bootstrap} --- the seven $E_8$ slices as an audit raster, the cascade
      bridge, the M\"obius bootstrap.
\item \texttt{tfpt\_4\_frontier} --- honest status of $\eta_B$, $m_p/m_e$, Koide, dark matter, full QG.
\item \texttt{tfpt\_5\_redteam} --- the adversarial audit: declared attacks, what survives, what each
      reduces to, and the kill tests.
\end{enumerate}
\noindent\textbf{Companions:} \textbf{R.} \texttt{tfpt\_research\_contracts} --- the open interfaces
$(v_{\mathrm{geo}}, G_{\mathrm{net}}, F_{\mathrm{transfer}})$ as numbered contracts;
\textbf{H.} \texttt{tfpt\_horizon\_readouts} --- Appendix~H (unit-system reframe): one seam constant as
the universal horizon thermal code, the Nariai/anchor surface, the resummed seam clock;
\textbf{O.} \texttt{origin\_theory} --- why no free fundamental number remains (exact core $+$ one
honestly-typed cyclic interpretation).
\end{navbox}

\noindent\textbf{You are reading \textbf{document~3} ($E_8$ Audit \& Bootstrap).}

% ---------------------------------------------------------------------
%  tfpt_status.tex -- THE canonical "status at a glance" card.
%  Single source for the overall theory status; \input near the front
%  of every document so there is ONE status statement, not one per
%  paper.  Requires tfpt_style.tex (keybox, markers).  Detailed,
%  paper-specific status expansions live only where they belong
%  (introduction: five-questions + closure ledger; tfpt_4: frontier
%  table; tfpt_research_contracts: the contracts).
% ---------------------------------------------------------------------
\begin{keyboxnb}[TFPT status at a glance --- read this card the same way in every document]
\textbf{Two inputs, one compiler.} From the boundary constant $c_3=\tfrac1{8\pi}$ (axiom~P1) and the
five-slot carrier $g_{\mathrm{car}}=5$ (axiom~P2) --- jointly the elementary symmetric data of one anchor
$a=(1,1,2)$ --- the discrete compiler $D_5\oplus A_3+\mu_4\Rightarrow E_8$ is \emph{built}, and the
Standard-Model skeleton (gauge group, three families, hypercharges, the flavor matrix), the dimensionless
constants ($\alpha^{-1}=137.0359992$) and the scale grammar are read off as fixed points. The algebraic
core is machine-checked; physical readouts run through \emph{named bridges} whose status is typed below.
\smallskip\\
\noindent\textbf{What is closed, and what genuinely remains.} The discrete/algebraic layer is closed
(\mE); the everything-open residual reduces to \emph{three named interface problems}:
\begin{center}\small
\renewcommand{\arraystretch}{1.18}
\begin{tabular}{@{}lll@{}}
\toprule
\textbf{Interface} & \textbf{Question} & \textbf{Status} \\
\midrule
$v_{\mathrm{geo}}$ & the one metrology unit ($=1/\sqrt G=m/\mu$); No-Unit Thm: no compiler scale & primitive \mO \\
$G_{\mathrm{net}}$ & simple-current extension $(D_5)_1{\otimes}(A_3)_1{\rtimes}\langle(1,1)\rangle\cong(E_8)_1$ & \mE/\mC \\
$F_{\mathrm{transfer}}$ & one functor, four typed interfaces (Koide, $\eta_B$, axion, $m_p/m_e$) & \mC \\
\bottomrule
\end{tabular}
\end{center}
\noindent Everything carries a claim ID resolving to a row of \texttt{verification/status\_ledger.csv}
(the single source of truth); \textbf{if the text and the ledger ever disagree, the ledger wins.} The
historical labels $(U_{\mathrm{wall}},G_{\mathrm{metric}},F_{\mathrm{frontier}})$ are retained only for
ledger continuity --- the live residual is $v_{\mathrm{geo}}\oplus G_{\mathrm{net}}\oplus F_{\mathrm{transfer}}$.
\markerkey
\end{keyboxnb}

\TFPTbadges
\TFPTclaimstack[compiler]


\tableofcontents
\clearpage
\part{The $E_8$ secondary-branching atlas}
\par\smallskip\noindent\textbf{Purpose and discipline.}
This note does \emph{not} introduce new physics. It builds an \emph{audit raster}: the claim that
\emph{every load-bearing TFPT number should appear in at least one $E_8$ branching projection}. A number
that appears in none is suspect; a number appearing in several independent projections becomes
structurally strong. All arithmetic identities below are exact \mE; the \emph{physical readings} are
audit-level \mO\ unless tied to a derivation. (All 77 numeric claims were machine-checked with exact
integer/rational arithmetic.)

% =====================================================================
\section{The compiler core (where this note sits)}
% =====================================================================
The strongest current form of the theory is a five-layer compiler:

\begin{center}
\begin{tikzpicture}[font=\footnotesize,
  L/.style={rectangle,rounded corners,draw,thick,align=center,inner sep=4pt,minimum height=8mm},
  ar/.style={-{Stealth[length=2mm]},thick,tfgray}]
  \node[L,draw=tfred!80!black,fill=tfred!8] (l1) at (0,0) {\textbf{L1 Inputs}\\$\cthree{=}\tfrac1{8\pi}$, $\gcar{=}5$ ($3{+}2$)};
  \node[L,draw=tfblue,fill=tfblue!8] (l2) at (3.3,0) {\textbf{L2 Geometry}\\$C^+{=}D_5$, $\PP^1\!\setminus\!\mu_4{=}A_3$\\$(D_5{\oplus}A_3){+}\mu_4{=}E_8$};
  \node[L,draw=tfblue,fill=tfblue!8] (l3) at (6.8,0) {\textbf{L3 Flavor}\\$E_\bullet{=}\mathcal O(-2){\oplus}\mathcal O(-1)^2$\\$a{=}(1,1,2)$, $R,L$};
  \node[L,draw=tfgold!80!black,fill=tfgold!12] (l4) at (10.1,0) {\textbf{L4 Seed}\\$\phiz,\ \alpha,\ \rho_\Lambda,\ M_{\mathrm{scal}}$};
  \node[L,draw=tfgreen!75!black,fill=tfgreen!8] (l5) at (12.9,0) {\textbf{L5 Readout}\\SM, CKM,\\PMNS, $A_s$};
  \draw[ar] (l1)--(l2); \draw[ar] (l2)--(l3); \draw[ar] (l3)--(l4); \draw[ar] (l4)--(l5);
\end{tikzpicture}
\end{center}

\noindent $E_8$ is now a \emph{closure} object, not an input: the carrier is the $D_5=\mathfrak{so}_{10}$
half-spinor, the four punctures $\PP^1\setminus\mu_4$ give $A_3=\mathfrak{su}_4$, and the common
discriminant group $\Z_4$ glues them, with the norm test $q(D_5)+q(A_3)=\tfrac54+\tfrac34=2$. This note
charts the \emph{secondary} maximal subalgebras of $E_8$ as audit windows onto the TFPT modules.

% =====================================================================
\section{The $(1,1,2)$ anchor microcode}
% =====================================================================
The parabolic anchor $a=(1,1,2)$ (splitting type of $E_\bullet=\mathcal O(-2)\oplus\mathcal O(-1)^2$) is a
miniature compiler. Its power sums $p_n(a)=1^n+1^n+2^n=2+2^n$ generate the discrete budget:

\begin{center}\small
\renewcommand{\arraystretch}{1.2}
\begin{tabularx}{\linewidth}{@{}clY@{}}
\toprule
$n$ & $p_n=2+2^n$ & \textbf{TFPT reading (verified)} \\
\midrule
1 & $4$ & $|\mu_4|=-\deg E_\bullet$ (glue index) \\
2 & $6$ & $|R^+(A_3)|$ (positive $A_3$ roots) $=\Z_6$ cusp order \\
3 & $10$ & $\AL=\binom52=\dim V_{D_5}$ (action ladder / pair sector) \\
4 & $18$ & $\Nfam\,|R^+(A_3)|=3\cdot6$ \\
5 & $34$ & $2h(D_5)+3|R^+(A_3)|=16+18$; and $34-|R(A_3)|=22=\sum R$ \\
\bottomrule
\end{tabularx}
\end{center}

\noindent and the budget compresses to products of power sums:
\[
  q(A_3)=\frac{p_2}{2p_1}=\frac34,\quad q(D_5)=\frac{p_3}{2p_1}=\frac54,\quad
  \textstyle\sum L=p_1p_3=40,
\]
\[
  \Oadm=2p_1p_2=48,\quad D_{\mathrm{start}}=p_2p_3=60,\quad |R(E_8)|=4p_2p_3=240.
\]
The first three power sums alone reach $E_8$: $p_1p_2p_3=4\cdot6\cdot10=240=|R(E_8)|$, and
$p_1p_2p_3+(p_4-p_3)=240+8=248=\dim E_8$.

\begin{winbox}[The anchor difference ladder: the binary carrier spine \mE]
While the \emph{products} of the power sums give the large budgets, the \emph{differences} give the
binary carrier spine in one line:
\[
  \boxed{\ p_n(a)=2+2^n\ \Longrightarrow\ \Delta p_n=p_{n+1}-p_n=2^n\ }:\quad
  (2,4,8,16)=\bigl(|\Z_2|,\,|\mu_4|,\,h(D_5),\,\dim S^+\bigr).
\]
So the single anchor $a=(1,1,2)$ generates the big numbers multiplicatively and the doubling spine
$2,4,8,16$ additively --- one vector, two readings.
\end{winbox}

\begin{winbox}[The anchor reconstructs the carrier rank --- and the $41$-chain]
The characteristic polynomial of the anchor vector is $(t-1)^2(t-2)=t^3-4t^2+5t-2$, with elementary
symmetric data
\[
  e_1(a)=4=|\mu_4|,\qquad e_2(a)=5=\gcar,\qquad e_3(a)=2=|\Z_2|_{\mathrm{sheet}} .
\]
So the parabolic flavor geometry reads the carrier rank $\gcar=5$ \emph{backwards}: $D_5\Rightarrow a$ but
also $a\Rightarrow\gcar$. Moreover $3^2+4^2=5^2$ sits as $(p_0,e_1,e_2)=(3,4,5)$, giving the
\emph{anchor form} of the abelian index
\[
  \boxed{\ 10\,b_1=e_1(a)^2+e_2(a)^2=4^2+5^2=41\ }
\]
alongside the known $41=40+1=\sum L+N_\Phi$. \mE
\end{winbox}

\begin{keybox}[Pascal closure: $\gcar=5$ is the unique solution (kills alternatives) \mE]
The carrier rank is not merely a good fit --- it is the \emph{unique} root of an exact closure. The
even-exterior carrier dimension $2^{g-1}$ must equal the truncated Pascal sum $\binom g0+\binom g1+\binom g2$
(the $1{+}g{+}\binom g2$ that builds $\dim S^+$):
\[
  \boxed{\ 2^{\,g-1}=\binom g0+\binom g1+\binom g2\ \Longleftrightarrow\ g=5\ }
\]
($g{=}3:4{\ne}7$; $g{=}4:8{\ne}11$; $\mathbf{g{=}5:16{=}16}$; $g{=}6:32{\ne}22$; $g{=}7:64{\ne}29$). At
$g=5$ the row is $\binom50,\binom51,\binom52=1,5,10$, the same Pascal triple that reads family, action
ladder and the $E_8$ root count $16\cdot5\cdot3=240$. This is the fourth independent lock on $\gcar=5$
(with rank-fill, Coxeter-match, and integer-glue/norm), and unlike those it explicitly \emph{kills} the
neighbours $g=4,6$.
\smallskip\\
\noindent\textbf{The Pascal Ladder Theorem (\veri{v108\_pascal\_ladder.py}).} The closure generalises
exactly: for \emph{every} truncation degree $K$,
\[
  \boxed{\ 2^{\,g-1}=\textstyle\sum_{k\le K}\binom gk\ \Longleftrightarrow\ g=2K+1\ }
\]
(odd $g$ solves by the binomial midpoint identity; even $g$ is excluded by the strict straddle
$\sum_{k<g/2}<2^{g-1}<\sum_{k\le g/2}$). \textbf{Consequence (zero slack): the carrier selection $g=5$ is
\emph{equivalent} to the truncation degree $K=2$} --- the Pascal-\emph{selection} residual of the red team
(Target~B) is therefore exactly the one named hypothesis \emph{Quadratic Boundary Locality} (the seam
certifies code degrees $\le2$; the architecture document). Neighbour worlds on the ladder: $K{=}1\to g{=}3$ (a
\emph{one}-family world), $K{=}2\to g{=}5$ (three families), $K{=}3\to g{=}7$ (\emph{nine} families),
$K{=}4\to g{=}9$ \emph{inconsistent} ($N_{\mathrm{fam}}=255/9\notin\Z$). Overdetermination: among the
ladder worlds, only $K=2$ also satisfies the independent rank closure $g+N_{\mathrm{fam}}=8=\operatorname{rank}E_8$
--- two independent selections agree on the same world. The hypothesis itself stays \mC.
\end{keybox}

\begin{winbox}[The seed is an anchor expression]
The retained seed is exactly the two-term anchor combination of $\cthree$:
\[
  \boxed{\ \phiz=\frac{2p_1}{p_2}\,\cthree+2p_1p_2\,\cthree^4=\frac43\cthree+48\,\cthree^4
  =\frac1{6\pi}+\frac{3}{256\pi^4}\ }
\]
(verified to $30$ digits). The Cabibbo base is $\lambda_Y=\sqrt{\phiz(1-\phiz)}$, the hierarchy engine
of the whole fermion spectrum. \mE
\end{winbox}

% =====================================================================
\section{The residue and length matrices as spectral objects}
% =====================================================================
With the residue matrix $R$ and the full word-length matrix $L=R+6W$ ($W$ the rank-one winding),
\[
  R=\begin{pmatrix}1&3&0\\1&5&2\\2&5&3\end{pmatrix},\qquad
  L=\begin{pmatrix}7&3&0\\7&5&2\\8&5&3\end{pmatrix},\qquad
  W=\begin{pmatrix}1&0&0\\1&0&0\\1&0&0\end{pmatrix},
\]
the spectral invariants are pure compiler numbers (all verified):
\begin{center}\small
\renewcommand{\arraystretch}{1.25}
\begin{tabularx}{\linewidth}{@{}lYlY@{}}
\toprule
\textbf{Invariant of $R$} & \textbf{value} & \textbf{Invariant of $L$} & \textbf{value} \\
\midrule
$\operatorname{tr}R$ & $9=\Nfam^2$ & $\operatorname{tr}L$ & $15$ \\
$\det R$ & $8=h(D_5)$ & $\det L$ & $20=|R(A_4)|=2\AL$ \\
$\chi_R(t)$ & $t^3-9t^2+10t-8$ & $\chi_L(t)$ & $t^3-15t^2+40t-20$ \\
$\operatorname{PrinMin}_2(R)$ & $(2,3,5)$ & $\operatorname{PrinMin}_2(L)$ & $(5,14,21)$ \\
$\|R\|_F^2$ & $78=\dim E_6$ & --- & --- \\
\bottomrule
\end{tabularx}
\end{center}

\noindent The three principal $2{\times}2$ minors of $R$ are $(2,3,5)$ --- sheet, family, carrier ---
with product $2\cdot3\cdot5=30=h(E_8)$ and sum $10=\AL$. The wrong down-branch $\{1,3,4\}$ fails exactly
these invariants.

\paragraph{Bilinear anchor forms.} With $\onev=(1,1,1)^\top$ and $a=(1,1,2)^\top$ (all verified):
\begin{center}\small
\renewcommand{\arraystretch}{1.2}
\begin{tabular}{@{}c|cc@{\hspace{2.5em}}c|cc@{}}
\toprule
$R$ & $\onev$ & $a$ & $L$ & $\onev$ & $a$ \\
\midrule
$\onev^\top$ & $22=\textstyle\sum R$ & $27=3^3$ & $\onev^\top$ & $40=|R(D_5)|$ & $45=\dim D_5$ \\
$a^\top$ & $32=2^5$ & $40=|R(D_5)|$ & $a^\top$ & $56=\dim\mathbf{56}_{E_7}$ & $64=\dim S^+\!\cdot|\mu_4|$ \\
\bottomrule
\end{tabular}
\end{center}

\noindent The corner value $a^\top L\,a=64=(16,4)$ reads exactly one spinor-glue sheet of the
$E_8$ branching. And the rank-one winding update shows up in the inverses:
\[
  \boxed{\ R^{-1}\onev=\tfrac14(1,1,-1)^\top,\qquad L^{-1}\onev=\tfrac1{10}(1,1,-1)^\top\ }
\]
i.e.\ the residue response reads the glue denominator $|\mu_4|=4$, the full-length response reads the
decuple denominator $\AL=10$. \mE

\begin{winbox}[Anchor-plane Plücker coordinates: scattered hits become oriented areas \mE]
The strongest anti-numerology tool is to read \emph{orientation invariants} of the anchor plane
$\langle\onev,a\rangle$ rather than loose integers. For $M\in\{R,Q,K,L\}$ the left block
$B_M=\bigl(\onev^\top M;\,a^\top M\bigr)$ (a $2{\times}3$) and the right block $C_M=(M\onev\ \ Ma)$ (a
$3{\times}2$) have the $2{\times}2$ Plücker minors
\[
  \renewcommand{\arraystretch}{1.15}
  \begin{array}{l|ccc}
   & \operatorname{Pl}(\cdot)\ \text{(left)} & \operatorname{Pl}_R(\cdot)\ \text{(right)} & \\\hline
  R & (-6,2,14) & (8,12,4)=(h(D_5),|R(A_3)|,|\mu_4|) & \\
  Q & (3,6,3) & (0,4,5)=(0,|\mu_4|,\gcar) & \\
  K & (-1,6,4)=(-N_\Phi,|R^+\!(A_3)|,|\mu_4|) & (12,12,-2)=(|R(A_3)|,|R(A_3)|,-|\Z_2|) & \\
  L & (6,26,14) & (20,30,10)=(2\AL,h(E_8),\AL) & \\
  \end{array}
\]
Three load-bearing flavor numbers are now single oriented areas, not coincidences:
\[
  \boxed{\ \|\operatorname{Pl}(K)\|_1=11,\qquad
  \|\operatorname{Pl}_R(K)\|_1=\operatorname{Pl}(L)_{13}=26,\qquad
  \sum\operatorname{Pl}_R(L)=60=D_{\mathrm{start}}\ }
\]
($11$ is the mass-power anchor-plane norm; $26$ the $c_t/c_b$ denominator; $60$ the $E_8$ cascade start).
The left-null-space response complements the right one above: the column sums of $\det L\cdot L^{-1}$ and
$\det R\cdot R^{-1}$ are $(-5,1,6)$ and $(1,-5,6)$, the same magnitudes
$\{N_\Phi,\gcar,|R^+\!(A_3)|\}=\{1,5,6\}$. Finally the sector rows themselves are anchor-shifted:
$K_e{-}K_u=(1,1,2)=a$ and $L_e{-}L_u=(1,2,3)=\operatorname{Exp}(A_3)$ --- leptons are up-quarks displaced by
the family roots. All entries are exact \mE; this is an audit layer, not a closure (the physical quark
amplitudes stay \mC). \veri{v37\_plucker\_anchor.py}
\end{winbox}

\begin{winbox}[The numerology null test: the look-elsewhere burden quantified \mE]
The raster above disciplines \emph{which} integers may appear; the complementary statistical question ---
\emph{could a random ``theory'' of equal formula complexity land this many data hits?} --- is answered by an
explicit, fully declared null model. For each of the $13$ scored frozen-registry observables the
\emph{entire} grammar class is enumerated exactly (no sampling): one-term expressions
$\tfrac pq\,(\phiz)^a\lamC^c(1{+}\lamC)^d\pi^b e^{u/w}$ with $p,q\le130$, $|a|\le2$, $|c|\le3$,
$|d|,|b|\le1$, $|u|\le6$, $w\le6$, and two-term sums with coefficients $\le12$ --- a class that provably
contains every TFPT formula it scores (machine-checked fairness condition). Data windows are conservative:
the \emph{maximum} of the achieved deviation, the documented experimental $\sigma$, the quote resolution and
the $5\%$ quark scheme spread, so the known tension observables ($\sin^2\theta_{13}$, $s_{23}$,
$\delta_{\mathrm{CKM}}$, $m_\mu/m_\tau$, $\beta$) enter with their honest \emph{large} windows and contribute
almost nothing. Result of the census: joint formula-fishing probability
$\prod_i p_i=10^{-25.8}$, stable to $<2$ orders of magnitude when the coefficient bound is varied
$80\dots200$ and the plausibility window $5{\times}\dots20{\times}$. Monte-Carlo cross-check
(deterministic seeds): $2\times2{\cdot}10^5$ pseudo-theories --- random rational dials inside TFPT's own
formula skeletons, with the seed $\phiz$ fixed resp.\ drawn log-uniformly --- reach at most $5/13$
resp.\ $4/13$ hits versus TFPT's $13/13$. The test has \emph{power}: perturbing $\phiz$ by
$1\%/10\%/30\%$ collapses TFPT's own scorecard to $9/6/4$ hits, and shuffling which measured value is
assigned to which formula collapses it to a mean of $1.2$. Separately, all $94\,500$ complexity-matched
variants of the $F_{U(1)}$ fixed-point equation (coefficients, $M$, transport exponent, resolvent power)
are root-solved: exactly \emph{one} --- the TFPT instance $(2,\tfrac45,M{=}41,e^{-2\alpha},-\tfrac54)$ ---
lands inside the achieved $4{\cdot}10^{-8}$ CODATA window, so $p_\alpha\le1.1{\cdot}10^{-5}$. Headline:
\[
  \boxed{\ P\bigl(\text{null reproduces the scorecard incl.\ }\alpha\bigr)
  \;\le\;10^{-30.7}\ \ (\approx102\ \text{bits})\ }
\]
\emph{conditional on the declared grammar} --- a look-elsewhere-corrected null-model rejection, never
``certainty''. Complementary to the frozen registry: the freeze (\vref{v84}) prevents \emph{future}
look-elsewhere drift, this test quantifies the \emph{retrospective} look-elsewhere burden.
\veri{v100\_numerology\_null\_mc.py}
\end{winbox}

% =====================================================================
\section{The secondary branching atlas}
% =====================================================================
Each maximal subalgebra of $E_8$ mirrors a TFPT module. The dimension identities are exact \mE; the
readings are audit-level \mO. The seven \emph{slices} of $\dim E_8=248$ are shown below: the same total
cut seven ways, each cut a different TFPT module.

\begin{figure}[h]\centering
\begin{tikzpicture}[x=0.0495cm,y=1cm,font=\scriptsize]
  % helper: \seg{xstart}{width}{ybase}{fillcolor}{label}
  \def\bh{0.62}
  \newcommand{\seg}[5]{%
    \fill[#4] (#1,#3) rectangle (#1+#2,#3+\bh);
    \draw[white,line width=0.4pt] (#1,#3) rectangle (#1+#2,#3+\bh);
    \node[font=\scriptsize\bfseries,text=white] at (#1+#2/2,#3+\bh/2) {#5};}
  % row y-levels
  \foreach \nm/\y in {D_8/6,D_5{\times}A_3/5,E_6{\times}A_2/4,E_7{\times}A_1/3,F_4{\times}G_2/2,A_4{\times}A_4/1,A_8/0}{
    \node[anchor=east,font=\footnotesize] at (-6,\y+\bh/2) {$\nm$};}
  % D8: 120,128
  \seg{0}{120}{6}{tfblue!75}{120}\seg{120}{128}{6}{tfgold!75!black}{128}
  % D5xA3: 45,15,60,64,64
  \seg{0}{45}{5}{tfblue!75}{45}\seg{45}{15}{5}{tfblue!45}{15}\seg{60}{60}{5}{tfgreen!60!black}{60}\seg{120}{64}{5}{tfgold!75!black}{64}\seg{184}{64}{5}{tfgold!60!black}{64}
  % E6xA2: 78,8,81,81
  \seg{0}{78}{4}{tfred!70!black}{78}\seg{78}{8}{4}{tfred!45!black}{8}\seg{86}{81}{4}{tfgreen!55!black}{81}\seg{167}{81}{4}{tfgreen!42!black}{81}
  % E7xA1: 133,3,112
  \seg{0}{133}{3}{tfblue!75}{133}\seg{133}{3}{3}{tfblue!40}{3}\seg{136}{112}{3}{tfgold!75!black}{112}
  % F4xG2: 52,14,182
  \seg{0}{52}{2}{tfblue!75}{52}\seg{52}{14}{2}{tfblue!45}{14}\seg{66}{182}{2}{tfgray!60}{182}
  % A4xA4: 24,24,50x4
  \seg{0}{24}{1}{tfblue!75}{24}\seg{24}{24}{1}{tfblue!55}{24}
  \foreach \i in {0,1,2,3}{\seg{48+\i*50}{50}{1}{tfgreen!55!black}{50}}
  % A8: 80,84,84
  \seg{0}{80}{0}{tfblue!75}{80}\seg{80}{84}{0}{tfgold!75!black}{84}\seg{164}{84}{0}{tfgold!60!black}{84}
  \draw[-{Stealth[length=1.6mm]},tfgray] (0,-0.5)--(248,-0.5) node[midway,below,font=\scriptsize]{$\dim E_8=248$};
\end{tikzpicture}
\caption{\small The $E_8$ slice atlas: seven maximal-subalgebra cuts of $248$. Gold $=$ spinor-glue
($64{=}(16,4)$, $112{=}7{\cdot}16$, $84{=}7{\cdot}12$, $128$); red $=$ flavor-residue ($78{=}\|R\|_F^2$,
$8{=}\det R$); green $=$ family/action blocks; blue $=$ Lie adjoints.}
\end{figure}

\begin{center}\small
\renewcommand{\arraystretch}{1.3}
\begin{tabularx}{\linewidth}{@{}p{2.1cm}p{4.1cm}Y@{}}
\toprule
\textbf{Subalgebra} & \textbf{$248=$} & \textbf{TFPT reading} \\
\midrule
$D_8$ & $120+128$ & family-transition adjoint $+$ two spinor-glue sheets ($128{=}2\dim S^+|\mu_4|$) \\
$D_5\times A_3$ & $45+15+60+64+64$ & \textbf{main split}: $\dim D_5$, $\dim A_3$, $\AL|R^+(A_3)|$, two $(16,4)$ glue sheets \\
$E_6\times A_2$ & $78+8+81+81$ & \textbf{flavor residue}: $\|R\|_F^2$, $\det R{=}\dim A_2$, cubic family blocks \\
$E_7\times A_1$ & $133+3+112$ & \textbf{scalaron}: $112{=}7\cdot16$ (exponent $\times$ one family) \\
$F_4\times G_2$ & $52+14+182$ & \textbf{occupancy/gauge}: $|R(F_4)|{=}48{=}\Oadm$, $|R(G_2)|{=}12{=}\dim\mathfrak g_{\mathrm{SM}}$ \\
$A_4\times A_4$ & $24+24+4\cdot50$ & \textbf{action ladder}: $24{+}24{=}48{=}\Oadm$; off-diagonal $(5,10)$ reps $\AH,\AL$ \\
$A_8$ & $80+84+84$ & rank-decuple $+$ scalaron: $80{=}8\cdot10$, $84{=}7\cdot12$ \\
\bottomrule
\end{tabularx}
\end{center}

\begin{theorem}[$E_6\times A_2$ flavor-residue shadow \mO]
The principal flavor branching reads the residue matrix:
\[
  \|R\|_F^2=78=\dim E_6,\qquad \det R=8=\dim A_2=h(D_5),\qquad \onev^\top R\,a=27,
\]
and the full $E_6\times A_2$ dimension count is the residue identity
\[
  \boxed{\ 248=\|R\|_F^2+\det R+2\,(\onev^\top R\,a)\,\Nfam=78+8+2\cdot27\cdot3 .\ }
\]
Here $E_6$ reads the Frobenius norm of the residue matrix, $A_2$ the pure three-family symmetry, and the
$\mathbf{27}$ the cubic family block. (Strongest new $E_8$ flavor fingerprint.)
\end{theorem}

\begin{theorem}[$E_7\times A_1$ scalaron shadow \mO]
With the scalaron exponent $7=\Oadm-10b_1=-\deg E+\operatorname{rk}E$,
\[
  248=133+3+112,\qquad 112=7\cdot\dim S^+=7\cdot16,\qquad |R(E_7)|=126=7\cdot18,
\]
where $18=\Nfam|R^+(A_3)|$. The scalaron exponent $7$ (the seam power $M_{\mathrm{scal}}^2/\bar M_{\mathrm{Pl}}^2=c_3^7$)
is thus exactly the $E_7$ off-block coefficient. \emph{Cross-check:} the centred $E_8$ exponents give
$\sum_{m\in\mathrm{Exp}(E_8)}|m-15|=56=\dim\mathbf{56}_{E_7}$ (exponents $\{1,7,11,13,17,19,23,29\}$,
the primitive residues mod $30$).
\end{theorem}

\begin{proposition}[$A_4\times A_4$ action-ladder and $F_4\times G_2$ occupancy \mO]
$248=24+24+4(5\cdot10)$ with $24+24=48=\Oadm$ and off-diagonal reps $(\mathbf 5,\mathbf{10})=(\AH,\AL)$
--- the action ladder appears inside $E_8$ (an audit, \emph{not} an $SU(5)$ GUT). And
$F_4\times G_2$: $|R(F_4)|=48=\Oadm$, $|R(G_2)|=12=\dim\mathfrak g_{\mathrm{SM}}$,
$\operatorname{rank}F_4=4=|\mu_4|$, $\operatorname{rank}G_2=2=|\Z_2|$.
\end{proposition}

\par\smallskip\noindent\textbf{Flavor-matrix and prime-table fingerprints (audit-level, machine-verified).}
Five exact linear-algebra readouts of $R,L$ and the prime atoms, kept here under the no-free-pattern rule.
\begin{itemize}\setlength\itemsep{2pt}
\item\textbf{Row norms ladder.} $\|R_u\|^2,\|R_d\|^2,\|R_e\|^2=(10,30,38)=(\AL,\,h(E_8),\,2\cdot19)$ with
$10{+}30{+}38=78=\dim E_6$ ($19\in\operatorname{Exp}(E_8)$).
\item\textbf{Residue inventory.} The entries of $R$ on $C_6$ are $m=(1,2,2,2,0,2)$: the \emph{only} gap is
$4=|\mu_4|=\sum a$. Moments $\sum R=22$, $\sum R^2=78=\dim E_6$ --- so $\|R\|_F^2=78$ is the second moment
of the $C_6$ inventory with the $\mu_4$ hole.
\item\textbf{First-column norm $\to$ $E_6{\times}A_2$ off-blocks.} $L_{\cdot1}=(7,7,8)$, $\|L_{\cdot1}\|^2=162=81+81$
--- the two $\mathbf{81}$ off-blocks of the $E_6{\times}A_2$ slice are exactly the wound first-generation column.
\item\textbf{Lepton exponent polynomial.} The $\phiz$-ladder lepton powers $(k_e,k_\mu,k_\tau)=(5,3,2)$ are
the Coxeter primes $30=2\cdot3\cdot5$ in hierarchy order: $(t-5)(t-3)(t-2)=t^3-10t^2+31t-30$
($e_1=\AL$, $e_3=h(E_8)$). The lepton $c$-ratios obey $\tfrac{(16/7)(7/6)}{4/3}=2=|\Z_2|$ and
$\tfrac{(16/7)(7/6)}{(4/3)^2}=\tfrac32$ (the near-Koide factor).
\item\textbf{The $2,3,5$ multiplication table} (the shortest integer-landscape compression):
\[
  (2,3,5)\xrightarrow{\times}(6,10,15,30)\xrightarrow{}(8{=}3{+}5,\,16,\,240{=}8\cdot30,\,248{=}240{+}8).
\]
\end{itemize}
\par\smallskip

% =====================================================================
\section{Group-theoretic audits (optional)}
% =====================================================================
\par\smallskip\noindent\textbf{Weyl completion index.}
Over the main stabiliser,
\[
  \frac{|W(E_8)|}{|W(D_5)|\,|W(A_3)|}=\frac{696729600}{1920\cdot24}=15120=63\cdot240=(2^6-1)\,|R(E_8)| .
\]
A group-theory machine-check, not physics: it quantifies the Weyl symmetry created beyond the
$D_5\times A_3$ stabiliser. \mE\ (arithmetic) / \mO\ (reading)
\par\smallskip

\par\smallskip\noindent\textbf{$E_8$ theta series as a higher-shell raster.}
$\Theta_{E_8}=E_4=1+240\sum_{n\ge1}\sigma_3(n)q^n$ has shell degeneracies
$240,\ 2160,\ 6720,\ 17520,\dots$ The first is $|R(E_8)|=240$; the second is
$2160=9\cdot240=\Nfam^2|R(E_8)|$. If later one-loop or CMB shell degeneracies touch these multiplicities,
the theta series is the right place to type them --- as an audit raster, not a claim. \mO
\par\smallskip

% =====================================================================
\section{Further audit lemmas (machine-verified this round)}
% =====================================================================
Five more exact identities tie the flavor matrices and the $E_8$ exponents to the compiler. All checked
with exact integer arithmetic.

\par\smallskip\noindent\textbf{(i) Full-length Frobenius norm is three $E_6$ blocks.}
$\|L\|_F^2=234=\Nfam\,\|R\|_F^2=3\dim E_6$, and
$\|L\|_F^2=\|R\|_F^2+12\langle R,W\rangle+36\|W\|_F^2=78+48+108$,
with $\langle R,W\rangle=4=|\mu_4|$ and $\|W\|_F^2=3=\Nfam$. The full word-length load is exactly the
residue $E_6$ norm $+$ the $A_3$-winding over $\mu_4$ $+$ the family-winding square. \mE

\par\smallskip\noindent\textbf{(ii) $A_3$-winding determinant lemma.}
The rank-one winding $L=R+6\,\mathbf 1 e_1^\top$ and the matrix-determinant lemma with
$e_1^\top R^{-1}\mathbf 1=\tfrac14$ give
$\det L=\det R\bigl(1+\tfrac{|R^+(A_3)|}{|\mu_4|}\bigr)=8(1+\tfrac64)=20=2\AL$.
So $h(D_5)=8\xrightarrow{\,A_3^+/\mu_4\,}20=|R(A_4)|$: the winding injects the carrier factor $5$ into
the cokernel ($\operatorname{coker}R\simeq\Z_8$, $\operatorname{coker}L\simeq\Z_{20}$). \mE

\par\smallskip\noindent\textbf{(iii) The scalaron $7$ is the $E_8$ exponent variance per decuple.}
Centring the exponents at $h/2=15$, $\sum_m(m-15)^2=560$ and
$\operatorname{Var}=\tfrac{560}{8}=70=7\cdot\AL=(\text{scalaron }7)\cdot(\text{decuple }10)$.
The halved positive offsets $\{1,2,4,7\}$ satisfy $\sum=14=2\cdot7$, $\sum{}^2=70$, $\prod=56=\dim\mathbf{56}_{E_7}$
--- a micro-code carrying the scalaron, the decuple and the $E_7$ minuscule at once.
Also $\sum_{\text{pairs}}m(30-m)=620=\tfrac{h\dim E_8}{12}=2\cdot10\cdot31$. \mE

\par\smallskip\noindent\textbf{(iv) Exponents $+$ degrees give the $120{+}128$ split.}
The invariant degrees $d_i=m_i+1=\{2,8,12,14,18,20,24,30\}$ obey
$\sum_i m_i=120$, $\sum_i d_i=128$, $248=120+128$, $\prod_i d_i=|W(E_8)|$,
matching the $D_8$ branching $248=120+128$; the degrees are a TFPT dictionary
($2,8,12,14,18,20,24,30$ $=$ sheet, $h(D_5)$, $|R(A_3)|$, $\dim G_2$, $p_4$, $\det L$, $\dim A_4$, $h(E_8)$). \mE

\par\smallskip\noindent\textbf{(v) Root shell $=$ local $E_7\times A_1$.}
Around any $E_8$ root the inner-product multiplicities are
$\{+2{:}1,\ +1{:}56,\ 0{:}126,\ -1{:}56,\ -2{:}1\}$, i.e.\ the local
$248=(133,1){+}(1,3){+}(56,2)$. A distinguished root (seam/Higgs direction) leaves an orthogonal $E_7$
and an active coupling shell $56{+}56$. \mE

% =====================================================================
\section{What this changes, and what stays open}
% =====================================================================
$E_8$ is one container whose maximal subalgebras each project a different TFPT module:

\begin{figure}[h]\centering
\begin{tikzpicture}[font=\footnotesize]
  \def\R{3.2}
  \node[circle,draw=tfgold!80!black,fill=tfgold!15,very thick,minimum size=15mm,align=center] (e8) at (0,0) {$E_8$\\{\scriptsize$248$}};
  \foreach \ang/\nid/\sub/\mod/\col in {
     90/na/{$D_5{\times}A_3$}/{SM carrier $+$ family glue}/tfblue,
     38/nb/{$E_6{\times}A_2$}/{flavor residue $\|R\|_F^2{=}78$}/tfred,
     -14/nc/{$E_7{\times}A_1$}/{scalaron $112{=}7{\cdot}16$}/tfgold,
     -62/nd/{$F_4{\times}G_2$}/{occupancy $48$ $+$ gauge $12$}/tfgreen,
     -118/ne/{$A_4{\times}A_4$}/{action ladder $5,10$}/tfgreen,
     -166/nf/{$A_8$}/{rank-decuple $80{=}8{\cdot}10$}/tfblue,
     -214/ng/{$D_8$}/{$120{+}128$ glue}/tfblue}{
     \node[rectangle,rounded corners,draw=\col,fill=\col!10,align=center,inner sep=2.5pt,font=\scriptsize] (\nid) at (\ang:\R) {\textbf{\sub}\\{\scriptsize\mod}};
     \draw[-{Stealth[length=1.8mm]},thick,tfgray!70] (e8)--(\nid);}
\end{tikzpicture}
\caption{\small $E_8$ as an audit container: each maximal subalgebra projects one TFPT module. The
discipline rule --- every load-bearing number appears in $\ge1$ projection --- makes the structure
falsifiable.}
\end{figure}

\par\smallskip\noindent\textbf{Net.}
The atlas turns isolated fingerprints (e.g.\ $\|R\|_F^2=78=\dim E_6$, the scalaron $7$) into
\emph{branching projections} of one container. The discipline rule --- ``every load-bearing number must
appear in $\ge1$ projection'' --- makes the structure falsifiable: a number in no projection is suspect.
This is a \emph{program}, not a proof of new physics: all readings are \mO. The two genuinely
load-bearing, derived results remain the $\mu_4$ glue ($E_8=(D_5\oplus A_3)+\mu_4$, \mE\ \veri{v1\_e8\_glue.py}) and the
spectral selector of the flavor matrix ($\det R=h(D_5)$, $\operatorname{PrinMin}_2=(2,3,5)$, \mE\ \veri{v4\_flavor\_matrix.py}).

\noindent The remaining flavor-geometry step (historically ``H2'') is sharpened in the parabolic note to:
\emph{the $D_4$-equivariant stable parabolic structure on $\mathcal O(-2)\oplus\mathcal O(-1)^2$ realises
the unique branch with $\det R=8$ and $\operatorname{PrinMin}_2(R)=(2,3,5)$} --- no longer ``find the
matrix'', but ``prove this one geometric realisation''. (Status: since \vref{v118}--\vref{v122} the dictionary
\emph{values} are exact theorems --- hexagon $=$ sign-twisted family spectrum, addresses pinned, margins
forced; what remains is exactly the realisation statement above, the selector establishment R4$'$.)

\clearpage
\part{The $E_8$ cascade bridge}
\noindent\textbf{The connection in one sentence.}
The old E$_8$ cascade (Paper V1.06: the nilpotent-orbit chain $D_n=60-2n$, $n=0\ldots26$) and the new
compiler/atlas are \emph{the same even-integer $E_8$ spine}, read three ways: the cascade \emph{orders}
the numbers $\{8,10,\dots,60\}$ as a scale ladder, the atlas \emph{reads} them as $E_8$ subalgebra
branchings, and the compiler \emph{generates} them from the $(1,1,2)$ anchor. The cascade \emph{starts}
at $D_{\mathrm{start}}=60=p_2p_3$ and \emph{ends} at $8=h(D_5)=\det R$ --- the flavor selector.

% =====================================================================
\section{The old cascade, recalled}
% =====================================================================
Paper V1.06 builds a deterministic scale ladder from the nilpotent orbits of $E_8$. With centralizer
dimension $D=248-\dim\mathcal O$ and the Hasse chain restricted to $\Delta D=2$, a unique maximal chain
runs from $A_4{+}A_1$ ($D_0=60$) down to the regular orbit ($D_{26}=8$):
\[
  \boxed{\ D_n=60-2n,\quad n=0,\dots,26;\qquad \varphi_{n+1}=\varphi_n e^{-\gamma(n)},\quad
  \gamma(n)=\lambda\,[\ln D_n-\ln D_{n+1}]\ }
\]
\noindent The cascade arithmetic (endpoints, exponent rungs, IR-tail product, variance) is
machine-checked.\veri{v5\_e8\_cascade.py}
with the single structural normalisation $\lambda=\gamma(0)/(\ln248-\ln60)=0.58770$, fixed by a
curvature-extremum on the chain (no fit). The $E$-windows of the two-loop flow anchor it: $\alpha_3=1/(8\pi)=\cthree$
(the $E_8$ window) and $\alpha_3=1/(6\pi)$ (the $E_6$ window, near $\phiz$).

% =====================================================================
\section{The spine: every rung is a current number}
% =====================================================================
Reading $D=60-2n$ against today's compiler atoms, the secondary-branching atlas and the $(1,1,2)$ anchor
power sums $p_n=2+2^n$, \textbf{25 of the 27 rungs are load-bearing numbers} (machine-checked):

\begin{center}\scriptsize
\renewcommand{\arraystretch}{1.04}
\begin{tabularx}{\linewidth}{@{}cc>{\raggedright\arraybackslash}p{4.9cm}cc>{\raggedright\arraybackslash}p{4.9cm}@{}}
\toprule
$n$ & $D$ & \textbf{current meaning} & $n$ & $D$ & \textbf{current meaning} \\
\midrule
0 & 60 & $D_{\mathrm{start}}=p_2p_3=\AL\,|R^+\!(A_3)|$ \textbf{(start)} & 13 & 34 & $p_5$ anchor power $=2\cdot17$ \\
1 & 58 & $2\cdot29$ (E$_8$ exponent) & 14 & 32 & $2^5=2\dim S^+$ \\
2 & 56 & $\dim\mathbf{56}_{E_7}=a^\top L\mathbf 1=\sum|m{-}15|$ & 15 & 30 & $h(E_8)=2{\cdot}3{\cdot}5$ \textbf{(Coxeter rung)} \\
3 & 54 & $2\cdot27$, $27=\mathbf 1^\top Ra=3^3$ & 16 & 28 & $4\cdot7$ (half of $56$) \\
4 & 52 & $\dim F_4=|R(D_5)|{+}|R(A_3)|$ & 17 & 26 & $2\cdot13$ (old hadronic block) \\
5 & 50 & $A_4{\times}A_4$ block $=5{\cdot}10$ & 18 & 24 & $\dim A_4$; $24{+}24{=}48$ \\
6 & 48 & $\Oadm=|R(F_4)|=2p_1p_2$ & 19 & 22 & $\sum R=\mathbf 1^\top R\mathbf 1$ (residue sum) \\
7 & 46 & $2\cdot23$ (E$_8$ exponent) & 20 & 20 & $\det L=|R(A_4)|=2\AL$ \\
8 & 44 & $4\cdot11$ (E$_8$ exponent) & 21 & 18 & $p_4=\Nfam|R^+\!(A_3)|$ \\
9 & 42 & $6\cdot7=|R^+\!(A_3)|\cdot7$ & 22 & 16 & $\dim S^+$ (one generation) \\
10 & 40 & $|R(D_5)|=\sum L=\mathbf 1^\top L\mathbf 1$ & 23 & 14 & $\dim G_2=\sum L_d=8{+}6$ \\
11 & 38 & $2\cdot19$ (E$_8$ exponent) & 24 & 12 & $|R(A_3)|=\dim\mathfrak g_{\mathrm{SM}}$ \\
12 & 36 & $6^2$ (old EW block) & 25 & 10 & $\AL=\binom52$ (old cosmo block) \\
   &    &                                & 26 & \phantom{0}8 & $h(D_5)=\det R=2p_1$ \textbf{(end: flavor selector)} \\
\bottomrule
\end{tabularx}
\end{center}

\noindent Three structural facts jump out (all \mE\ arithmetic):
\begin{itemize}[leftmargin=1.4em]
\item \textbf{Endpoints are the compiler boundary.} The cascade \emph{starts} at $60=p_2p_3=D_{\mathrm{start}}$
      and \emph{ends} at $8=h(D_5)=\det R$ --- the flavor spectral selector is literally the deepest
      orbit of $E_8$.
\item \textbf{Coxeter rung is $h(E_8)$.} At $n=15$, $D=30=h(E_8)$ --- this is the \emph{Coxeter rung},
      not the arithmetic centre; the median node is $n=13$, $D=34=p_5(a)$ (below). It matches the
      exponent-centering audit
      $\sum_{m}|m-15|=56=\dim\mathbf{56}_{E_7}$ (which is itself rung $n=2$).
\item \textbf{Doubled $E_8$ exponents are rungs.} $\{58,46,38,34,26,22,14\}=2\times\{29,23,19,17,13,11,7\}$
      are exactly the cascade rungs at $n=1,7,11,13,17,19,23$; three of them coincide with anchor numbers
      ($34=p_5$, $22=\sum R$, $14=\dim G_2$). Only $n=8$ ($D=44=4\cdot11$) is not independently labelled.
\end{itemize}

\begin{winbox}[Cascade arithmetic --- endpoints, centre and length (verified)]
The $27$ rungs encode $E_8$ in their endpoints, and their centre is the anchor:
\begin{align*}
  &\tfrac{D_{\mathrm{start}}\,D_{\mathrm{end}}}{2}=\tfrac{60\cdot8}{2}=240=|R(E_8)|,\qquad
  240+D_{\mathrm{end}}=248=\dim E_8,\qquad D_{\mathrm{start}}=2\,h(E_8)=60;\\
  &\text{arithmetic median}=\tfrac{60+8}{2}=34=p_5(a)=h(E_8)+|\mu_4|,\qquad
  \textstyle\sum_{n=0}^{26}(D_n-30)=27\cdot4=27|\mu_4|;\\
  &\text{width}=D_{\mathrm{start}}-D_{\mathrm{end}}=52=\dim F_4,\qquad
  60\!\to\!30=h(E_8),\quad 30\!\to\!8=\textstyle\sum R=22;\\
  &\#\text{nodes}=27=3^3=\mathbf 1^\top R\,a\ (E_6\text{ cubic block}),\qquad
  \#\text{steps}=26=\dim\mathbf{26}_{F_4}.
\end{align*}
So the cascade is an $E_8$ Coxeter ladder centred on $p_5(a)$ with a per-node $|\mu_4|$ glue shift; its
length ($27$ nodes, $26$ steps) carries the $E_6$ cube and the $F_4$ fundamental. \mE
\end{winbox}

\begin{keybox}[The IR tail $n=24\ldots30$: the compiler atoms]
Continuing $D=60-2n$ past the last orbit ($D{=}8$, $n{=}26$) is no longer a nilpotent-orbit chain but the
\emph{algebraic compiler tail} of the Coxeter clock:
\[
  D_{24\ldots30}=\underbrace{12}_{|R(A_3)|=\dim\mathfrak g_{\mathrm{SM}}},\
  \underbrace{10}_{\AL},\ \underbrace{8}_{h(D_5)},\
  \underbrace{6}_{|R^+(A_3)|},\ \underbrace{4}_{|\mu_4|},\ \underbrace{2}_{\text{sheet}},\ \underbrace{0}_{\text{closure}} .
\]
This is exactly what the old ``$n=0\ldots30$ clock'' was scanning: $n{=}0\ldots26$ the $E_8$ orbit spine,
$n{=}27\ldots30$ the compiler atoms. (Mark the tail as the \emph{compiler tail}, not orbits.) \mE
\end{keybox}

\begin{winbox}[Three further cascade lemmas (all machine-verified)]
The clock carries more $E_8$ data than the endpoints:
\begin{itemize}\setlength\itemsep{3pt}
\item\textbf{Primitive rungs sum to the root count.} Evaluating $D_n=60-2n$ at the \emph{primitive}
indices $n\in(\Z/30)^\times=\{1,7,11,13,17,19,23,29\}=\mathrm{Exp}(E_8)$ gives
$\{58,46,38,34,26,22,14,2\}$ with
\[
  \sum_{n\in\mathrm{Exp}(E_8)}(60-2n)=240=|R(E_8)| ,
\]
and since $30-n$ is primitive when $n$ is, the clock indexes its own $E_8$ exponents.
\item\textbf{The IR tail product is the main Weyl stabiliser.} The nonzero tail
$\{12,10,8,6,4,2\}$ has
\[
  \prod=46080=|W(D_5)|\,|W(A_3)|=1920\cdot24,\qquad \textstyle\sum=42=|R^+(A_3)|\cdot7 .
\]
\item\textbf{Cascade variance $=$ $E_6$ flavour norm times the scalaron-gauge block.} The $27$ rungs
$60,\dots,8$ have mean $34=p_5(a)$ and
\[
  \sum_{n=0}^{26}(D_n-34)^2=6552=78\cdot84=\dim E_6\cdot(7\,\dim\mathfrak g_{\mathrm{SM}}) ,
\]
tying the cascade simultaneously to the $E_6{\times}A_2$ flavour shadow and the $A_8$ scalaron-gauge
shadow. \mE\ (audit level)
\end{itemize}
\end{winbox}

% =====================================================================
\section{Three views of one spine}
% =====================================================================
\begin{center}
\begin{tikzpicture}[font=\footnotesize,yscale=0.27,xscale=1]
  % vertical spine D=60..8
  \draw[very thick,gray!55] (0,8)--(0,60);
  \foreach \D in {8,10,12,16,18,20,24,30,40,48,52,56,60}{\fill[gray!55] (0,\D) circle (1.6pt);}
  % zone brackets
  \node[tfblue,anchor=west] at (0.35,58) {$\mathbf{60,56,52,50,48}$ --- E-series / atlas branchings (UV)};
  \node[tfgold!80!black,anchor=west] at (0.35,40) {$\mathbf{40}=|R(D_5)|=\sum L$};
  \node[tfred!80!black,anchor=west] at (0.35,30) {$\mathbf{30}=h(E_8)$ \,(Coxeter rung, $n{=}15$)};
  \node[tfgold!80!black,anchor=west] at (0.35,24) {$\mathbf{24,20,18}$ --- $\det L$, $p_4$, $A_4$};
  \node[tfgreen!60!black,anchor=west] at (0.35,15) {$\mathbf{16,12,10}=\dim S^+,|R(A_3)|,\AL$ (IR atoms)};
  \node[tfgreen!60!black,anchor=west] at (0.35,8) {$\mathbf{8}=h(D_5)=\det R$ \,(flavor selector, end)};
  \node[anchor=south] at (0,60.5) {\scriptsize $D{=}60$ start};
  \node[anchor=north] at (0,7.5) {\scriptsize $D{=}8$ end};
  % three readings
  \node[align=center,tfblue] at (8.6,52) {\textbf{cascade}\\\emph{orders} as $\varphi_n$\\scale ladder};
  \node[align=center,tfgold!80!black] at (8.6,32) {\textbf{atlas}\\\emph{reads} as $E_8$\\branchings};
  \node[align=center,tfgreen!60!black] at (8.6,14) {\textbf{compiler}\\\emph{generates} from\\$(1,1,2)$ anchor};
\end{tikzpicture}
\end{center}

\par\smallskip\noindent\textbf{The unification.}
The same numbers $\{8,10,\dots,60\}$ are: (a) the \emph{scale ladder} $\varphi_n$ of the old cascade
(orbit dimensions, ordered by $e^{-\gamma}$); (b) the \emph{branching blocks} of the secondary-atlas
($\dim F_4{=}52$, $\dim\mathbf{56}_{E_7}$, $5{\cdot}10$, \dots); (c) the \emph{anchor power sums and
matrix invariants} of the compiler ($p_n=2+2^n$, $\det R{=}8$, $\sum L{=}40$). The cascade is the
\emph{UV$\to$IR ordering}; the compiler is its \emph{IR endpoint}, where the small numbers
$8,10,12,16$ are the flavor/family/carrier data.

\begin{keybox}[The spine tetrahedron: the microkernel as \emph{one} finite object \mE\ (\veri{v91\_spine\_tetrahedron.py})]
The spine $T=\{2,3,4,5\}=\{|\Z_2|,\Nfam,|\mu_4|,\gcar\}$ equals $\{e_3(a),p_0(a),e_1(a),e_2(a)\}$ for the
anchor $a=(1,1,2)$ --- the tetrahedron on these four nodes is the anchor plus its zeroth moment, not a
second structure. Its combinatorics \emph{are} the central integer grammar: edge products
$\{6,8,10,12,15,20\}$ ($|R^+(A_3)|$; $\operatorname{rank}E_8{=}h(D_5){=}\det R$; $\AL{=}|E(K_5)|$;
$|R(A_3)|{=}\dim\mathfrak g_{\mathrm{SM}}$; $\dim A_3$; $\det L$), face products $\{24,30,40,60\}$
($|W(A_3)|$; $h(E_8)$; $|R(D_5)|{=}\sum L$; $D_{\mathrm{start}}$), volume
$2{\cdot}3{\cdot}4{\cdot}5=120=|R^+(E_8)|=5!$; and the graph reading
$|R(E_8)|=240=|\mu_4|\cdot|E(K_4)|\cdot|E(K_5)|$ with $|E(K_4)|=p_2=6$, $|E(K_5)|=p_3=10$ --- including
the honest negative control that the identification \emph{breaks} at $K_6$ ($|E(K_6)|=15\neq p_4=18$),
so it is specific, not generic. Two disciplined caveats: the ``dual cuts'' (node $\times$ opposite face
$=120$, etc.)\ are tautologies of any complementary $4$-set partition and are typed as presentation
only; and the tetrahedron grammar is a \emph{sub}-grammar --- the load-bearing $7,16,41,48,240,248$ are
\emph{not} subset products and still require the anchor power sums / Pascal / glue layers.
\end{keybox}

% =====================================================================
\section{Direct $\phiz$ relations: old vs current}
% =====================================================================
The old cascade's ``fundamental relations near $n=0$'' connect directly to the current note --- one is
\emph{identical}, two are early forms of current refinements:

\begin{center}\small
\renewcommand{\arraystretch}{1.25}
\begin{tabularx}{\linewidth}{@{}p{2.4cm}YYc@{}}
\toprule
\textbf{Quantity} & \textbf{old (V1.06)} & \textbf{current} & \textbf{relation} \\
\midrule
$\Omega_b$ & $\phiz(1-2\cthree)=0.04894$ & $(1-\tfrac1{4\pi})\phiz=0.04894$ & \textbf{identical} ($2\cthree=\tfrac1{4\pi}$) \mE \\
$V_{us}/V_{ud}$ & $\sqrt{\phiz}=0.2306$ & $\lambda_C=\sqrt{\phiz(1-\phiz)}=0.2244$ & refined (the $(1-\phiz)$ factor) \mE \\
$r$ (tensor) & $(\phiz)^2=0.00283$ & $12/N_\star^2=0.00397$ & two branches, same order \mC \\
\bottomrule
\end{tabularx}
\end{center}
\markerkey

\noindent So $\Omega_b=\phiz(1-2\cthree)$ was already the exact current formula in V1.06; the Cabibbo
ratio gained its $(1-\phiz)$ factor; and the old direct $r=(\phiz)^2$ and the current Starobinsky
$r=12/N_\star^2$ agree in order ($\sim0.003$) but are distinct predictions --- a clean, testable
discriminant, stated honestly.

% =====================================================================
\section{Sector inventory: have all the old-cascade sectors been computed?}\label{sec:sectorcov}
% =====================================================================
The old cascade also carried \emph{calibrated absolute scales}: per block, one unit $\zeta_B$ was fixed on
a single reference and every other scale in the block followed fit-free from the ratio law
$\varphi_m/\varphi_n=(D_m/D_n)^\lambda$ (Paper~V1.06, App.~C). These block assignments are the place where
the proton mass, the EW masses and the CMB temperature lived. The honest question is: \emph{has the new
compiler re-derived each of these, or only some?} The full inventory, with the current status spelled out:

\begin{center}\scriptsize
\renewcommand{\arraystretch}{1.25}
\begin{tabularx}{\linewidth}{@{}p{1.4cm}p{3.5cm}Yc@{}}
\toprule
\textbf{Old block (rung)} & \textbf{Old observable \& calibration} & \textbf{Current-model status} & \textbf{Mark} \\
\midrule
near $n{=}0$ & $\Omega_b=\phiz(1{-}2\cthree)$, $V_{us}/V_{ud}=\sqrt{\phiz}$, $r=(\phiz)^2$ & \textbf{covered} --- dimensionless, derived from $\{\cthree,\phiz\}$; see the $\phiz$-table above. $\Omega_b$ is \emph{identical}. & \mE \\
EW, $n{=}12$ & $v_H=\zeta_{\rm EW}M_{\mathrm{Pl}}\varphi_{12}$, $M_W{=}\tfrac12 g_2 v_H$, $M_Z$ & \textbf{scale covered, absolute scheme-level} --- the EW scale enters the current note via the action-grammar exponent and the two-loop flow; $M_W,M_Z$ are dimensionful and remain RG/scheme outputs (one $\zeta_{\rm EW}$), \emph{not} new compiler powers. & \mC \\
hadronic, $n{=}15,17$ & $\mathbf{m_p}{=}\zeta_p M_{\mathrm{Pl}}\varphi_{15}$, $m_b{=}\zeta_b M_{\mathrm{Pl}}\varphi_{15}$, $m_u{=}\zeta_u M_{\mathrm{Pl}}\varphi_{17}$ & \textbf{ratios closed, absolute scale an anchor.} The quark mass \emph{ratios} sit on the $\phiz$-word-length ladder and are closed (integer Pl\"ucker readouts, \vref{v49}/\vref{v71}); the \emph{absolute} values need the per-sector $\Lambda$/$U_f^\star$ normalisation, an anchor (\mO). The \textbf{proton mass} is QCD-confinement, $m_p\sim\Lambda_{\rm QCD}$ by dimensional transmutation --- a \emph{cross-sector} ratio, listed as a genuine frontier item ($m_p/m_e$, tfpt\_4). & \mO \\
cosmo, $n{=}25$ & $T_{\gamma0}{=}\zeta_\gamma M_{\mathrm{Pl}}\varphi_{25}$, $T_\nu{=}(\tfrac4{11})^{1/3}T_{\gamma0}$, $\rho_\Lambda$ & \textbf{partly.} $\rho_\Lambda$/$\Lambda$ is \emph{typed} by the action grammar and the seam ($\Lambda$-typing, tfpt\_2); $T_\nu/T_{\gamma0}$ is standard QFT (not cascade). The \textbf{absolute CMB temperature} $T_{\gamma0}$ is a downstream readout of the reheating pipeline (the old cosmology/CMB layers; today the scalaron-reheating chain, \vref{v86}), conditional --- not re-derived as a compiler number. & \mC \\
\bottomrule
\end{tabularx}
\end{center}
\markerkey

\begin{gapbox}[Direct answer: which sectors are closed, which are not]
\textbf{Closed (derived from $\{\cthree,\gcar,\phiz\}$):} the dimensionless near-$n{=}0$ anchors
($\Omega_b$ exactly, Cabibbo, tensor branch) and all the \emph{dimension/rung} structure of the spine.
\textbf{Scale-level only (one $\zeta_B$ each, scheme/RG):} the EW masses $M_W,M_Z$ and the cosmological
absolute scales $T_{\gamma0}$ --- exactly as in the current SM/cosmology notes; these were \emph{never}
pure predictions, not even in V1.06 (each needed a block unit). \textbf{Genuinely open:} the
\textbf{proton mass} $m_p$ (and $m_p/m_e$) is QCD-confinement across two sectors and is \emph{not} claimed
as a compiler power --- it stays a frontier item (\mO, tfpt\_4); the quark mass \emph{ratios} are closed
(integer Pl\"ucker, \vref{v49}/\vref{v71}) and only the \emph{absolute} quark scale $m_b,m_u$ needs the per-sector
$\Lambda$/$U_f^\star$ normalisation, an anchor (\mO). So: yes, the \emph{spine} is fully audited and the
\emph{dimensionless} sectors are derived; the \emph{dimensionful} hadronic/cosmo absolute scales are
honestly still scheme-level, with the proton mass the one explicitly open cross-sector number --- nothing
is silently claimed.
\end{gapbox}

% =====================================================================
\section{What is new, what is old, what stays open}
% =====================================================================
\begin{keybox}[New structures found via the bridge]
\begin{enumerate}[leftmargin=1.5em]
\item The flavor selector $\det R=h(D_5)=8$ is the \emph{terminal orbit} of the $E_8$ cascade; the
      compiler is the cascade's IR endpoint. \mE
\item The atlas branching blocks ($52,56,60$) are the cascade's \emph{UV head} ($n=0,2,4$); the secondary
      atlas and the orbit cascade are the same spine read top-down vs.\ as subalgebras. \mE
\item The anchor power sums sit on doubled-exponent rungs: $p_5=34$ ($n{=}13$), $\sum R=22$ ($n{=}19$),
      $\dim G_2=14$ ($n{=}23$). \mE
\item The cascade Coxeter rung $D=30=h(E_8)$ ($n{=}15$) ties to the exponent-centering audit
      $\sum|m{-}15|=56$. \mE
\end{enumerate}
\end{keybox}

\begin{gapbox}[Honest separation]
The \emph{dimension/rung coincidences} are exact arithmetic \mE\ --- they show the old and new
constructions live on one $E_8$ spine. The old cascade's \emph{scale assignments} $\varphi_n$ (EW at
$n{=}12$, hadronic at $n{=}15,17$, cosmo at $n{=}25$) are a \emph{different} mechanism from the current
$\phiz$-power mass ladder: the cascade orders \emph{orders of magnitude} (which sector sits where) via
$e^{-\gamma}$, while the $\phiz$-ladder gives \emph{values within} a sector. Both are real, complementary
orderings; neither is forced onto the other. The per-sector $\zeta$-calibrated absolute scales (proton
mass, EW masses, CMB temperature) remain scheme-level --- audited one by one in
\S\ref{sec:sectorcov}, with the proton mass the single explicitly open cross-sector number.
\end{gapbox}

\noindent\textbf{Bottom line.} The old E$_8$ cascade was an early, correct intuition: the physically
relevant numbers are exactly the even-integer $E_8$ orbit spine $\{8,\dots,60\}$. The current work
\emph{explains why} --- they are the $(D_5\times A_3)+\mu_4$ compiler atoms and their $E_8$ branching
blocks --- and identifies the spine's two ends as the compiler boundary ($D_{\mathrm{start}}=60$,
$\det R=8$). Old and new are the same structure at different resolution.

\clearpage
\part{The self-consistency loop}
\noindent\textbf{The question, and the honest answer.}
The theory rests on two inputs, $\cthree=\tfrac1{8\pi}$ and $\gcar=5$. The old theory had a Möbius seam
and the seed was the \emph{start} of the $E_8$ cascade. Run it \emph{backwards}: can the cascade/$E_8$
\emph{return} $\cthree$ and $\gcar$, closing a self-consistent loop? \textbf{Answer:} the \emph{discrete}
content is a closed, overdetermined fixed point --- $\gcar=5$ and the ``$8$'' in $\cthree$ are forced by
the $E_8$ closure. The \emph{only} irreducible primitive left is the geometric $\pi$ (the Möbius/boundary
Gauss--Bonnet). So ``two axioms'' collapses to ``\emph{one} continuous primitive $\pi$ $+$ \emph{one}
discrete self-consistency fixed point.''

% =====================================================================
\section{The loop for $\gcar$ --- an overdetermined fixed point}
% =====================================================================
Two \emph{independent} structural conditions, both pure Lie-theory, force $\gcar=5$ once $E_8$ is the
closure object and $A_3$ is the four-puncture family geometry:

\begin{keybox}[$\gcar=5$ is forced twice over]
\begin{description}[leftmargin=2.2em,style=nextline]
  \item[\textnormal{(A) rank fill.}] The glue $D_{\gcar}\oplus A_3\hookrightarrow E_8$ must fill the
  rank: $\operatorname{rank}D_{\gcar}+\operatorname{rank}A_3=\operatorname{rank}E_8$, i.e.\
  $\gcar+3=8\Rightarrow\gcar=5$.
  \item[\textnormal{(B) Coxeter match.}] The carrier Coxeter number must equal the $E_8$ rank,
  $h(D_{\gcar})=2\gcar-2=8\Rightarrow\gcar=5$.
\end{description}
Both give $\gcar=5$, and the loop \emph{closes}: $\gcar=5\Rightarrow D_5\oplus A_3\xrightarrow{\mu_4}E_8
\Rightarrow h(E_8)=30=2\cdot3\cdot5\Rightarrow$ largest prime $=5=\gcar$. \mE
\end{keybox}

\begin{keybox}[Reverse glue selection --- $D_5$ \emph{and} $A_3$ from unimodular closure]
The closure even \emph{selects the two factors}, not just the rank. The discriminant of $A_{\mu-1}$ is
$\Z_\mu$ and of $D_g$ (odd $g$) is $\Z_4$; a common cyclic glue forces $\mu=4$, hence $A_{\mu-1}=A_3$.
The even-glue norm $q(D_g)+q(A_3)=\tfrac g4+\tfrac34=2$ then forces $g=5$, hence $D_5$. A second,
independent confirmation of $\mu=4$ is the \emph{mirror identity}
$\dim D_{\mu+1}+\dim A_{\mu-1}=\dim V_{D_{\mu+1}}\cdot|R^+(A_{\mu-1})|$, i.e.\ $3\mu(\mu+1)=\mu(\mu+1)(\mu-1)
\Rightarrow\mu-1=3\Rightarrow\mu=4$ (the $45+15=10\cdot6$ identity). So
\[
  \boxed{\ \text{common }\Z_4\text{ glue }+\text{ root norm }2\ \Longrightarrow\ D_5+A_3\ }
\]
--- $\gcar=5$ is recovered from the closure, not just posited. \mE
\end{keybox}

\begin{keybox}[Familyful $E_8$-closure fixed-point theorem]
\emph{Among rank-filling pairs $D_g\oplus A_m$ ($g+m=8$) admitting an integer glue index
$\sqrt{\det D_g\,\det A_m}=\sqrt{4(m+1)}$, exactly two close to an even unimodular lattice: $D_8\oplus A_0$
(no family) and $D_5\oplus A_3$ (familyful). The unique familyful solution is $D_5\oplus A_3$, with
lattice $\cong E_8$ and $\gcar=5$.}
\begin{center}\scriptsize\renewcommand{\arraystretch}{1.05}
\begin{tabular}{@{}lcccl@{}}
\toprule
$D_g{+}A_m$ & $g{+}m$ & glue idx$^2{=}4(m{+}1)$ & integer? & status \\
\midrule
$D_8{+}A_0$ & $8$ & $4$ & $2$ & valid, \emph{no family} \\
$D_7{+}A_1$ & $8$ & $8$ & --- & no integer glue \\
$D_6{+}A_2$ & $8$ & $12$ & --- & no integer glue \\
$\mathbf{D_5{+}A_3}$ & $8$ & $16$ & $\mathbf 4{=}|\mu_4|$ & \textbf{valid, familyful, $E_8$} \\
$D_4{+}A_4$ & $8$ & $20$ & --- & no integer glue \\
\bottomrule
\end{tabular}
\end{center}
Equivalently, with $q(A_{\mu-1})=\tfrac{\mu-1}{\mu}$ and $g=9-\mu$ the norm closure
$q(D_g)+q(A_{\mu-1})=2$ becomes $\mu^2-5\mu+4=0$, whose nontrivial root is $\mu=4$ ($A_3$, $g=5$). So
$\gcar=5$ is forced \emph{three} ways: rank-fill, Coxeter-match, and integer-glue/norm closure. \mE\ \veri{v6\_bootstrap.py}
\end{keybox}

\begin{keybox}[Bootstrap classification: $D_5\oplus A_3$ is the \emph{unique} familyful $E_8$ glue \mE]
$E_8$ is the unique even unimodular rank-$8$ lattice. Its two-factor root-lattice gluings $L_1\oplus L_2$
($\operatorname{rank}=8$) that admit a \emph{cyclic} glue are exactly those whose discriminant groups are
isomorphic and cyclic:
\[
  D_5\oplus A_3\ (\Z_4),\quad E_6\oplus A_2\ (\Z_3),\quad E_7\oplus A_1\ (\Z_2),\quad A_4\oplus A_4\ (\Z_5).
\]
Imposing the two TFPT data --- one factor with a $16$-dimensional half-spinor ($\dim S^+_{D_g}=2^{g-1}=16\Leftrightarrow g=5$)
and one factor giving exactly $\Nfam=3$ families ($\operatorname{rank}A_m=m=3$) --- selects
\[
  \boxed{\ D_5\oplus A_3\ \text{(uniquely), glue group }\Z_4=|\mu_4|\ } .
\]
$E_6{\oplus}A_2$ gives only $2$ families, $E_7{\oplus}A_1$ only $1$; $D_8$ has non-cyclic glue
$(\Z_2)^2$ and $A_8$ is a single rank-$8$ factor. So the carrier$+$family decomposition is forced.
\veri{v15\_bootstrap\_classification.py}
\end{keybox}

\begin{figure}[h]\centering
\begin{tikzpicture}[font=\small,
  n/.style={rectangle,rounded corners,draw,thick,align=center,inner sep=4pt,minimum height=8mm},
  ar/.style={-{Stealth[length=2.4mm]},very thick,tfgold!70!black}]
  \def\R{3.0}
  \node[n,draw=tfblue,fill=tfblue!8] (g) at (90:\R) {$\gcar=5$\\(carrier rank)};
  \node[n,draw=tfblue,fill=tfblue!8] (d5) at (162:\R) {$D_5=\mathfrak{so}_{10}$\\$h(D_5)=8$};
  \node[n,draw=tfgreen!70!black,fill=tfgreen!8] (e8) at (234:\R) {$(D_5{\oplus}A_3){+}\mu_4$\\$=E_8$, rank $8$};
  \node[n,draw=tfgold!80!black,fill=tfgold!12] (h) at (306:\R) {$h(E_8)=30$\\$=2\cdot3\cdot5$};
  \node[n,draw=tfblue,fill=tfblue!8] (p) at (18:\R) {max prime\\$=5$};
  \draw[ar] (g) to[bend right=12] (d5);
  \draw[ar] (d5) to[bend right=12] (e8);
  \draw[ar] (e8) to[bend right=12] (h);
  \draw[ar] (h) to[bend right=12] (p);
  \draw[ar] (p) to[bend right=12] (g);
  \node[tfgray] at (0,0) {\footnotesize self-consistent};
\end{tikzpicture}
\caption{\small The closed loop: $\gcar=5$ builds $E_8$, whose Coxeter number $30=2\cdot3\cdot5$ returns
the carrier prime $5$. The value $\gcar=5$ is the unique fixed point (overdetermined by rank-fill
\emph{and} Coxeter-match).}
\end{figure}

\noindent The carrier split $3+2$ and the one-generation count $\dim S^+=2^{\gcar-1}=16$ are then
automatic, as is the glue norm $q(D_5)+q(A_3)=\tfrac54+\tfrac34=2$. The family rank is
$\operatorname{rank}A_3=3=\Nfam$ and the glue index is $h(A_3)=4=|\mu_4|$ --- the Coxeter numbers of the
three players are $\bigl(h(A_3),h(D_5),h(E_8)\bigr)=(4,8,30)=\bigl(|\mu_4|,\ \det R,\ \text{compiler}\bigr)$.

% =====================================================================
\section{The loop for $\cthree$ --- the ``$8$'' is the $E_8$ rank}
% =====================================================================
The boundary normalisation is $\cthree=1/(8\pi)$. Its integer is exactly the carrier Coxeter number,
which the loop above equals to the $E_8$ rank:
\[
  \boxed{\ \cthree=\frac{1}{h(D_5)\,\pi}=\frac{1}{(2\gcar-2)\,\pi}=\frac{1}{8\pi},\qquad
  8=h(D_5)=\operatorname{rank}E_8=\det R=\varphi(30)\ }
\]
where $\varphi(30)=8$ is Euler's totient (the number of primitive residues mod $30$ $=$ the $E_8$
exponents) --- so the ``$8$'' has \emph{five} concordant readings. The tree seed
$\varphi_{\mathrm{tree}}=1/(6\pi)$ has $6=|R^+(A_3)|=h(E_8)/\gcar$, the three Gauss--Bonnet boundary
cycles ($6\pi=3\times2\pi$). So the \emph{discrete} parts of both $\cthree$ and $\varphi_{\mathrm{tree}}$
are $E_8$/carrier invariants; only the factor $\pi$ is genuinely continuous.

\begin{gapbox}[Honest caveat on $8\pi$]
$8\pi$ is \emph{also} the standard Einstein normalisation ($G_{\mu\nu}=8\pi G\,T_{\mu\nu}$). The loop
\emph{identifies} this gravitational $8$ with $\operatorname{rank}E_8=h(D_5)$ --- a non-coincidence
\emph{if} $E_8$ is the closure, but one should be clear that $8\pi$ has an independent GR reading. The
loop promotes the two $8$'s to the same number; it does not derive $\pi$.
\end{gapbox}

% =====================================================================
\section{The Möbius origin of the irreducible $\pi$}
% =====================================================================
What the loop does \emph{not} remove is $\pi$. In the old theory this is exactly the Möbius datum: the
orientable double cover of the Möbius fibre is a cylinder with two boundaries plus one $\Z_2$ seam
$\Gamma$; Gauss--Bonnet gives $\oint k\,ds=2\pi+2\pi+2\pi=6\pi$ (hence $\varphi_{\mathrm{tree}}=1/(6\pi)$),
and the $\Z_2$ identification \emph{is} the sheet ``$2$'' of $30=2\cdot3\cdot5$. So the Möbius seam
supplies precisely the two things the discrete loop cannot: the continuous $\pi$ and the sheet parity
$\Z_2$.

\par\smallskip\noindent\textbf{Net: two axioms $\to$ one primitive $+$ one fixed point.}
\begin{keyeq}
$\displaystyle \{\cthree,\gcar\}\ \longrightarrow\
  \underbrace{\pi\ (\text{Möbius/Gauss--Bonnet})}_{\text{one continuous primitive}}\ +\
  \underbrace{E_8\text{ closure}}_{\text{one discrete fixed point}}$
\end{keyeq}
\noindent The discrete data $\{2,3,4,5,6,8,30\}$ are \emph{all} self-consistent outputs of the $E_8$ closure
($\gcar=5$, $h(D_5)=8$, $h(A_3)=4=|\mu_4|$, $h(E_8)=30=2{\cdot}3{\cdot}5$, $|R^+(A_3)|=6$). The only
irreducible input on the \emph{dimensionless} axis is $\pi$. \mE\ (discrete loop) / \mO\ ($\pi$ is primitive)
\par\smallskip
\noindent On the \emph{dimensionful} axis there is exactly \emph{one} more irreducible: a single mass scale.
But even that is no longer a \emph{free} constant --- fixing the physical $\Lambda$ branch pins $\bar M_{\rm Pl}$,
so $G_N$ is a $\Lambda$-metrology \emph{output} (\texttt{origin\_theory}, $\Lambda$-typing); the one
genuinely irreducible scale is the induced-gravity (Sakharov) anchor. So the complete honest reduction is
``$\pi$ $+$ the $E_8$-closure fixed point $+$ \emph{one} dimensionful scale''.

\noindent\emph{One sharpening of the ``$8$''.} Both $\pi$-coefficients are $2\pi$ times a compiler integer:
$\varphi_{\mathrm{tree}}=1/(6\pi)$ with $6=2\Nfam$ (three Gauss--Bonnet boundary circles) and
$\cthree=1/(8\pi)$ with $8=2|\mu_4|$ (seam winding), so $\varphi_{\mathrm{tree}}=(|\mu_4|/\Nfam)\cthree=\tfrac43\cthree$.
The \emph{same} integer $8$ is the Hawking/Einstein coefficient ($G_{\mu\nu}=8\pi T_{\mu\nu}$, Appendix~H):
the seam normaliser, the $E_8$ rank and the universal horizon factor are one number, overdetermined
($8=2|\mu_4|=\operatorname{rank}E_8=h(D_5)=\varphi(30)=\det R$). \mE\ \veri{v54\_seam\_horizon\_keystones.py}

\begin{keybox}[A cyclic element \emph{inside} the hull: the order-$30$ Coxeter rotation \mE\ (\veri{v55\_coxeter\_cycle.py})]
The compiler hull carries a \emph{literal} cyclic symmetry. Built from the $E_8$ Cartan matrix, the Coxeter
element $c=s_1\cdots s_8$ has \textbf{order exactly $h(E_8)=30=|\Z_2|\cdot\Nfam\cdot\gcar=2\!\cdot\!3\!\cdot\!5$}
(verified by direct matrix powers), and its eigenvalues are the primitive $30$th roots
$e^{2\pi i\,m/30}$ with $m$ the $E_8$ \emph{exponents} $\{1,7,11,13,17,19,23,29\}$ --- which are precisely the
$\varphi(30)=8$ \emph{totatives} of $30$. Hence
\[
  \boxed{\ \operatorname{rank}E_8=8=\varphi(30)=\#\{\text{exponents}\}=\#\{\text{live phases of the order-}30\text{ cycle}\}\ }
\]
So the $8$ in $\cthree=\tfrac1{8\pi}$ is the count of live phases of a primitive order-$30$ cyclic rotation of
the hull, whose period $30=|\Z_2|\cdot\Nfam\cdot\gcar$ is the product of the three core integers. The cyclic
structure is \emph{present in the mathematics}, not interposed.
\end{keybox}

\begin{keybox}[The compiler atoms \emph{are} the $E_8$ Casimir degrees \mE\,(\veri{v66\_e8\_casimir\_degrees.py})]
The load-bearing integers are not an ad-hoc list: they are the fundamental \emph{invariant degrees} of $E_8$
(exponents $+1$),
\[
  \deg(E_8)=\{2,\ 8,\ 12,\ 14,\ 18,\ 20,\ 24,\ 30\},
\]
with the clean (non-trivial) readings $2{=}|\Z_2|$, $8{=}\operatorname{rank}E_8$, $14{=}\dim G_2$ (the $[K,Q]$
commutator $G_2$, flavor sector), $20{=}\det L$, $24{=}|W(A_3)|{=}4!$, $30{=}h(E_8)$ (and the fittable
$12{=}|R(A_3)|$, $18{=}p_4(a){=}2{+}2^4$). Two exact sums close it:
\[
  \begin{gathered}
  \textstyle\sum\deg(E_8)=128=2^{7}\quad(\text{the de Sitter prefactor }128\cthree^4;\ 7{=}\text{scalaron exp}),\\[2pt]
  \textstyle\sum\text{exp}(E_8)=120=|R^+(E_8)| .
  \end{gathered}
\]
Since $E_8$ is the audit hull, its invariant degrees being the theory's integers is \emph{structurally
expected} --- an organizing simplification that confirms the $E_8$ choice rather than a coincidence.
\end{keybox}

\begin{keybox}[Two-level foundation, and the parameter audit]
\textbf{Local foundation:} $\cthree,\gcar$ are the two operational inputs (the two axioms).
\textbf{Global bootstrap foundation:} $\gcar$ and the integer in $\cthree$ are $E_8$-closure fixed
points; only $\pi$ is primitive. The audit:
\begin{center}\small\renewcommand{\arraystretch}{1.2}
\begin{tabularx}{\linewidth}{@{}p{2.3cm}p{4.1cm}Y@{}}
\toprule
\textbf{former input} & \textbf{status now} & \textbf{recovered by} \\
\midrule
$\gcar=5$ & discrete fixed point & rank-fill, Coxeter-match, integer-glue/norm, $h(E_8)$ \\
$8$ in $\cthree$ & discrete fixed point & $h(D_5)$, $\operatorname{rank}E_8$, $\det R$, $\varphi(30)$, $2|\mu_4|$ (seam winding) \\
$6$ in $\varphi_{\mathrm{tree}}$ & discrete fixed point & $|R^+(A_3)|=h(E_8)/\gcar$ \\
$\pi$ & geometric primitive & Möbius / Gauss--Bonnet (\emph{not} derived) \\
$\phiz$ & derived readout & $\tfrac43\cthree+48\cthree^4$ \\
\bottomrule
\end{tabularx}
\end{center}
\textbf{Correct one-liner:} \emph{TFPT has no freely tunable fundamental numbers left; the only remaining
continuous origin is $\pi$ from the Möbius boundary geometry.} (Not ``zero axioms''.) \mE/\mO
\end{keybox}

% =====================================================================
\section{Honest status: bootstrap, not creation from nothing}
% =====================================================================
\begin{itemize}[leftmargin=1.4em]
\item \textbf{What is genuinely shown \mE:} $\gcar=5$ is an \emph{overdetermined} fixed point of the
  $E_8$ closure (two independent Lie-theory conditions); the integer ``$8$'' in $\cthree$ is the carrier
  Coxeter $=$ $E_8$ rank; the discrete data closes a loop. The earlier landscape scan confirms
  $\gcar=5$ is the \emph{unique} value at which all hierarchies/family/gauge close --- so the fixed point
  is unique, not one of many. \textbf{This uniqueness is now machine-formalised \mE\ in Lean 4} (Theorem A,
  \texttt{lean4-carrier-rigidity}): the division-free Pascal condition $2^{g}=g^2+g+2$ has the \emph{unique}
  solution $g=5$ (\texttt{carrier\_rank\_pascal\_unique}; growth lemma for $g\ge6$; $\Nfam=(2^{4}{-}1)/5=3$,
  $g{+}\Nfam=8$), audited to use only the standard kernel axioms. The 2026-06-11 hardening round adds two
  formal modules: \texttt{AnchorLadder.lean} (the anchor power-sum law $p_n(a)=2+2^n$ for \emph{all} $n$,
  the ladder atoms $(3,4,6,10)$, $240/8/248$, and the binary ladder $p_{n+1}{-}p_n=2^n$ as general
  theorems) and \texttt{PascalLadder.lean} (the odd-midpoint ladder identity
  $\sum_{k\le K}\binom{2K+1}k=2^{2K}$ for \emph{all} $K$ via Mathlib, plus the kernel-decidable bounded
  uniqueness of the $K{=}2$ carrier case) --- the solves-side of the Pascal Ladder Theorem
  (\veri{v108\_pascal\_ladder.py}) is thereby \mE.
\item \textbf{What this is \emph{not}:} a derivation from zero inputs. A self-consistency loop is
  consistent \emph{at} a value; uniqueness is what makes it predictive (and that is supplied by the
  landscape scan). The continuous $\pi$ is not produced by the loop --- it is the Möbius/boundary
  geometry, the one irreducible primitive.
\item \textbf{Conceptual gain:} the theory's foundation tightens from ``two free axioms'' to ``\emph{the}
  $E_8$-closure fixed point $+$ the geometric $\pi$.'' That is the Möbius self-consistency you intuited:
  the seed is the start of the cascade, and the cascade returns the seed's discrete content.
\end{itemize}

\clearpage
\part{Open items and the $E_8$-slice scan}

% =====================================================================
\section{The honest open-items list}
% =====================================================================
This is the complete current ledger of what is \emph{not} fully closed, graded by how hard the gap is.

\begin{center}\footnotesize
\renewcommand{\arraystretch}{1.25}
\begin{tabularx}{\linewidth}{@{}p{0.5cm}p{4.5cm}cY@{}}
\toprule
\# & \textbf{Open item} & \textbf{Grade} & \textbf{Hardening path / what closes it} \\
\midrule
\multicolumn{4}{@{}l}{\textit{1. SM core --- finite computations, no new physics}}\\
1a & absolute quark $c$-scale (ratios closed) & \mO & quark \emph{ratios} are integer Pl\"ucker readouts on the derived selector stratum (\vref{v49}/\vref{v69}/\vref{v71}); only the \emph{absolute} $U_f^\star$ normalisation remains --- an anchor \\
1b & $\theta_{12}$ solar angle & \mE & \textbf{texture closed (below):} explicit $\mu\tau$-symmetric texture verified; only $\varepsilon{=}\tfrac34\phiz{\approx}\cthree$ ($0.23\%$) stays conditional \\
\midrule
\multicolumn{4}{@{}l}{\textit{2. Scheme layer --- standard RG, no new input}}\\
2a & $m_W,m_Z,m_H,\sin^2\theta_W,\alpha_s$ & \mC & RG threshold matching $R_{\mathrm{SM}}$; $m_H$ on the closed UV quartic $1/16\pi^2$ \\
2b & hadron pole spectrum (singlet) & \mC & singlet algebra $\mathcal A_{\mathrm{sing}}$ explicit; dynamics already closed (gap $6\log\tfrac32$) \\
\midrule
\multicolumn{4}{@{}l}{\textit{3. Conditional --- named analytic hypotheses}}\\
3a & gravity (Hodge/$R^2$ branch) & \mC & APS-Fredholm $+$ truncation written out; full spectral UV completion remains \\
3b & strong CP / QFT closure & \mC & reflection positivity (RP) is the one remaining input \\
3c & $\Lambda$ branch label (rec vs car) & \mC & one discrete choice; typed via $(8\pi)^2$, value not fixed \\
\midrule
\multicolumn{4}{@{}l}{\textit{4. Axioms --- not eliminable, only hardenable}}\\
4 & P1 ($\cthree$), P2 ($\gcar{=}5$) & \mO & more Lean; P2 algebra already machine-verified; remain inputs \\
\midrule
\multicolumn{4}{@{}l}{\textit{5. Frontier --- honestly not on a clean ladder (see frontier doc)}}\\
5 & $\eta_B$, $m_p/m_e$, Koide, dark matter, full QG & \mO & genuine handles only; not forced (scan results below) \\
\midrule
\multicolumn{4}{@{}l}{\textit{6. Cosmology late-time pipeline --- not simplified by the compiler}}\\
6 & $C_\ell$ transfer, $a_{\ell m}$ seam, baryogenesis & \mC & Paper-7 machinery; $a_{\ell m}$ ``good CMB world, not yet this CMB world'' (SMICA test) \\
\bottomrule
\end{tabularx}
\end{center}
\markerkey

\noindent\textbf{Net.} Items 1 are finite arithmetic; 2 are standard RG; 3 hinge on named hypotheses
(RP, the $\Lambda$ branch, the UV completion); 4 are the two irreducible axioms; 5--6 are honestly open.
The inflation sector (the only cosmology piece the compiler \emph{does} sharpen) is now a parameter-free
prediction ($M{=}\cthree^{7/2}\Mbar$, document~2).

\begin{winbox}[Item 1b closed: the explicit solar-angle Majorana texture \mE]
The last open SM angle is now realised by an \emph{explicit} texture, not a fitted number. Take the
$\mu\tau$-symmetric Majorana matrix (tri-bimaximal at $\varepsilon{=}0$)
\begin{equation}
  M_\nu(\varepsilon)=\begin{pmatrix}-\eta(\varepsilon)&1&1\\ 1&1&0\\ 1&0&1\end{pmatrix},\qquad
  \eta(\varepsilon)=\frac{2\sqrt2}{\tan2\theta_{12}}-1,\quad
  \sin^2\theta_{12}=\tfrac13(1-2\varepsilon).
\end{equation}
$\mu\tau$-symmetry forces $\theta_{23}=45^\circ$ and $\theta_{13}=0$ \emph{exactly}; diagonalisation at the
TFPT solar deviation $\varepsilon=\tfrac34\phiz$ gives
\begin{equation}
  \boxed{\ \sin^2\theta_{12}=\tfrac13-\tfrac{\phiz}{2}=0.306747,\quad
  \theta_{23}=45^\circ,\quad \theta_{13}=0\ }\quad\mE
\end{equation}
\textbf{Honest residual} (\veri{v9\_neutrino\_texture.py}). The texture and the angle are exact. The seam identification is
$\varepsilon=\tfrac34\phiz=\tfrac34\bigl(\tfrac1{6\pi}+\tfrac3{256\pi^4}\bigr)=\cthree+\tfrac{9}{1024\pi^4}$:
the \emph{leading geometric term is exactly} $\cthree=\tfrac1{8\pi}$, and the $0.23\%$ residual is purely the
$\phiz$ seed tail (document~2). So $\varepsilon=\cthree$ holds at leading order; the full
$\varepsilon=\tfrac34\phiz$ is what enters the texture. The reactor angle $\theta_{13}$
($\sin^2\theta_{13}=e^{-5/6}\phiz$, document~2) enters through a separate $\mu\tau$-breaking term and is not
part of this minimal solar texture. \textbf{Scope:} this realises the TFPT solar shift in the
$\mu\tau$-symmetric \emph{limit}; it is \emph{not} the full PMNS matrix. The $\theta_{23}{=}45^\circ$ value
is the symmetric-limit centre, \emph{not} an octant prediction --- current global fits (NuFIT~6.0) leave
the octant unresolved --- and the reactor angle plus the combined matrix require the separately stated
$\mu\tau$-breaking channel.
\end{winbox}

% =====================================================================
\section{The $E_8$-slice scan: results (honest hit / miss)}
% =====================================================================
All seven maximal slices of $248$ were scanned for the unassigned quantities. The rule was a meaningful
(low-complexity) compiler/slice expression matching to $<0.5\%$; matches that need a free fitted factor
are reported as \emph{misses}.

\begin{center}\footnotesize
\renewcommand{\arraystretch}{1.25}
\begin{tabularx}{\linewidth}{@{}p{3.0cm}Yc@{}}
\toprule
\textbf{Slice / target} & \textbf{Result} & \textbf{Verdict} \\
\midrule
$A_4{\times}A_4=SU(5)^2$ & $\sin^2\theta_W|_{\mathrm{GUT}}=3/8$ sits in the $SU(5)$ hypercharge embedding; the $12$ coset states are the proton-decay $X,Y$ bosons & \mE\ (known) \\
$F_4{\times}G_2$ / Koide & Koide $Q=2/3=|\Z_2|/\Nfam$ (sheet over families); $J_3(\mathbb O)$ has $\dim 27{=}26{+}1$, $3{\times}3$ octonionic ``generations'' on the diagonal & \mC\ structural \\
$D_8=SO(16)$ & $128=(16,4)\oplus(\overline{16},\overline4)$ --- \emph{same} spinor as $D_5{\times}A_3$; $\nu_R$ (seesaw) lives in the $16$; \textbf{no new dark-matter state} (only the 4th $\mu_4$ direction) & \mE \\
$A_8=SU(9)$, $9{=}3{\times}3$ & $\Lambda^3(3,3)=(10,1){\oplus}(8,8){\oplus}(1,10)$: the $(8,8){=}64$ dominates --- naive family$\times$colour \textbf{fails} & miss \\
$A_8\supset SU(5){\times}SU(4)$ & $\Lambda^3(5{+}4)$ contains the SM $\mathbf{10}$: a genuine chiral-family-unification \emph{lead}, but not TFPT-forced & lead \\
$E_6{\times}A_2$ / Jarlskog & cross term $2(\mathbf1^\top Ra)\Nfam=162=248{-}78{-}8$ exact; $J\sim\lamY^6$ only to right \emph{order}, no clean hit & weak \\
$E_8$ theta shells & root shells $240,2160,6720,\dots$; $31{=}2^{\gcar}{-}1$ is \emph{not} a clean shell ratio & inconcl. \\
$m_p/m_e=1836.15$ & no clean $E_8$/compiler expression (a QCD-confinement ratio); apparent hits need a free factor & \textbf{miss} \\
$\eta_B\approx6.1{\times}10^{-10}$ & sits \emph{between} $\cthree^6{=}4{\times}10^{-9}$ and $\cthree^7{=}1.6{\times}10^{-10}$; right order, no integer power & order only \\
\bottomrule
\end{tabularx}
\end{center}
\markerkey

\begin{winbox}[Genuine new readings the scan produced]
\begin{itemize}\setlength\itemsep{2pt}
\item\textbf{Koide $=|\Z_2|/\Nfam$.} The empirical Koide value $2/3$ is exactly the sheet over the family
count --- a clean (if simple) compiler reading, promoting Koide from ``near $2/3$'' to ``democratic
$|\Z_2|/\Nfam$.''
\item\textbf{$E_8$ as the optimal $8$-bit code.} $E_8$ is the optimal $8$D lattice (kissing number $240$).
In byte framing $256{-}248=8=\operatorname{rank}E_8$ and $256{-}240=16=\dim S^+$ --- a clean
information-theoretic reading of the two load-bearing $E_8$ numbers.
\item\textbf{Dark matter $=$ no new state.} The $SO(16)$ scan shows the $128$ is the \emph{same} matter
spinor; there is no spare $E_8$ singlet to be a WIMP. Any DM candidate must be the determinant-line
axion (frontier doc), not a new $E_8$ rep --- a useful \emph{negative}.
\item\textbf{$SU(9)\supset SU(5){\times}SU(4)$ chiral lead.} $\Lambda^3$ puts the SM $\mathbf{10}$ in the
antisymmetric; the cleanest unexploited route to a family-unification reading, flagged for future work.
\end{itemize}
\end{winbox}

\begin{gapbox}[Honest negatives]
$m_p/m_e$ has \emph{no} clean $E_8$/compiler derivation (it is a confinement ratio --- the apparent
``hits'' all need a free fitted factor); $\eta_B$ matches only at the order-of-magnitude level
($\sim\cthree^{6.x}$); the $\Lambda$ rec/car branch is \emph{not} selected by the theta series; the
Jarlskog invariant is right only to order $\lamY^6$. These are reported as open, not forced.
\end{gapbox}

% =====================================================================
\section{Following the $SU(9)$ lead: cross-slice consistency}
% =====================================================================
The one scan lead worth following is $A_8=SU(9)\supset SU(5)\times SU(4)$. Splitting $9=(5,1)\oplus(1,4)$
the $E_8$ branching $248=80\oplus84\oplus\overline{84}$ decomposes \emph{exactly} as:

\begin{center}\small
\renewcommand{\arraystretch}{1.2}
\begin{tabularx}{\linewidth}{@{}llY@{}}
\toprule
piece & $SU(5)\times SU(4)$ content & reading \\
\midrule
$80$ (adj) & $(24,1){+}(1,15){+}(5,\bar4){+}(\bar5,4){+}(1,1)$ & $SU(5)_{\rm GUT}$ $+$ $SU(4)_{\rm fam}$ gauge \\
$84=\Lambda^3 9$ & $(10,1){+}(10,4){+}(5,6){+}(1,4)$ & the SM $\mathbf{10}$ as family-singlet $+$ family-$\mathbf4$ \\
$\overline{84}$ & $(\overline{10},1){+}(\overline{10},\bar4){+}(\bar5,6){+}(1,\bar4)$ & the conjugate matter \\
\bottomrule
\end{tabularx}
\end{center}

\noindent (Dimension checks: $24{+}15{+}20{+}20{+}1=80$; $10{+}40{+}30{+}4=84$; total $248$.)

\begin{winbox}[The genuine content: $A_8$ and $D_5{\times}A_3$ agree on the family \mE]
The $SU(4)$ here is \emph{the same} $A_3$ family group as in the construction slice $D_5\times A_3$ --- the
four punctures $\mu_4$ of $\PP^1\setminus\mu_4$. The $SU(5)$ is the Georgi--Glashow GUT sitting inside the
carrier $SO(10)=D_5$ ($SU(5)\subset SO(10)$). So the two maximal slices are \emph{consistent}: both read
the family as $SU(4){=}A_3$, and $A_8$ additionally exposes the $SU(5)$ matter content --- the SM
$\mathbf{10}$ appears as $(10,1)\oplus(10,4)$, i.e.\ a family-singlet $\mathbf{10}$ plus a family-$\mathbf4$,
where the $\mathbf4$ of $A_3$ is exactly TFPT's ``$4=\Nfam+1$'' (three families $+$ the $\mu_4$ direction).
\end{winbox}

\begin{gapbox}[Honest limit: chirality is still the index, not $E_8$]
$84\oplus\overline{84}$ is \emph{vectorlike} ($10\oplus\overline{10}$ etc.): $E_8$ being non-chiral, the
slice cannot by itself produce three \emph{chiral} generations. The chirality is the boundary index
$\operatorname{Ind}(D^+_{\rm fam})=3$ --- the \emph{same} input TFPT already uses. So $SU(9)$ is a
consistent re-reading (matter in $SU(5)$, family in $A_3$), not an independent derivation of $\Nfam=3$.
Each $SU(5)$ pair $(\mathbf{10}+\bar{\mathbf5})$ is anomaly-free, so the embedding is clean. \mC
\end{gapbox}

% =====================================================================
\section*{Appendix: computational verification}
\addcontentsline{toc}{section}{Appendix: computational verification}
% =====================================================================
The audit/bridge identities of this document (marked \mE, \mE) are re-derived from $\{\cthree,\gcar\}$ and
machine-checked by the suite in the repository folder \texttt{verification/}; the inline tags
\texttt{[verification/$\dots$]} point to the exact script. Relevant here:
\begin{itemize}[leftmargin=1.4em,itemsep=1pt]
\item \texttt{v1\_e8\_glue.py} --- the $\mu_4$ glue $E_8{=}(D_5{\oplus}A_3){+}\mu_4$ (load-bearing).
\item \texttt{v4\_flavor\_matrix.py} --- the spectral selector $\det R{=}h(D_5){=}8$, minors $\{2,3,5\}$.
\item \texttt{v5\_e8\_cascade.py} --- the cascade $D_n{=}60{-}2n$: endpoints, exponent rungs, IR tail, variance.
\item \texttt{v6\_bootstrap.py} --- the self-consistency loop: $\mu^2{-}5\mu{+}4{=}0$, $\gcar{=}5$ three ways, $8{=}\operatorname{rank}E_8$.
\end{itemize}
The $E_8$-slice/atlas readings are explicitly \mO\ (audit-level), not part of the pass/fail suite. Run
\texttt{python verification/run\_all.py}; full map in \texttt{verification/README.md}.

\end{document}
